% 手征对称性与杨--米尔斯梯度流（中文译本）
% 原文：M. Lüscher, JHEP 04 (2013) 123 — arXiv:1302.5246
% 逐节忠实翻译自英文转排版；公式照抄，参考文献保留英文。
\documentclass[11pt]{ctexart}
\usepackage[a4paper,margin=2.5cm]{geometry}
\usepackage{amsmath,amssymb,mathtools}
\usepackage{graphicx}
\usepackage{microtype}
\usepackage{enumitem}
\usepackage{url}
\usepackage[hidelinks]{hyperref}
\setCJKmainfont{AR PL UMing CN}
\setCJKsansfont{Droid Sans Fallback}
\setCJKmonofont{Droid Sans Fallback}

% ------------------------------------------------------------------
% 兼容层：原 plain-TeX 源的排版机制（与英文版一致）
% ------------------------------------------------------------------
\makeatletter
\newcount\equationno  \equationno=0
\def\equationlabel{}

\def\lussec #1. #2 \par{%
  \par\vskip 1.8em plus .4em minus .2em
  \gdef\equationlabel{#1}\global\equationno=0
  {\noindent\large\bfseries #1.\quad #2\par}\nobreak
  \vskip .9em plus .2em\noindent\ignorespaces}

\def\lussubsec #1 \par{%
  \par\vskip 1.5em plus .3em minus .2em
  {\noindent\normalsize\bfseries #1\par}\nobreak
  \vskip .7em plus .2em\noindent\ignorespaces}

\def\lusapp #1. #2 \par{%
  \par\vskip 1.8em plus .4em minus .2em
  \gdef\equationlabel{\hbox{\rm#1}}\global\equationno=0
  {\noindent\large\bfseries 附录 #1.\quad #2\par}\nobreak
  \vskip .9em plus .2em\noindent\ignorespaces}

\def\enum{\global\advance\equationno by \@ne
  \hfill\llap{{(\equationlabel.\the\equationno)}}}
\def\noenum{\hfill}
\def\equation#1{$$\displaylines{\qquad #1}$$}
\def\next#1{\cr\noalign{\vskip#1}\qquad}
\let\nexteq\next

% plain-TeX 展示对齐宏（依 plain.tex）
\def\openup#1{\advance\lineskip#1\advance\baselineskip#1\advance\lineskiplimit#1}
\def\ialign{\everycr{}\tabskip\z@skip\halign}
\def\displ@y{\mathsurround\z@\displayindent\z@\everycr{}\relax}
\def\@lign{\everycr{}}
\def\displaylines#1{\displ@y\tabskip\z@skip
  \halign{\hbox to\displaywidth{$\@lign\hfil\displaystyle##\hfil$}\crcr
    #1\crcr}}
\def\normalbaselines{\baselineskip\normalbaselineskip
  \lineskip\normallineskip\lineskiplimit\normallineskiplimit}
\def\cases#1{\left\{\vcenter{\normalbaselines\m@th
  \ialign{$##\hfil$&\quad##\hfil\crcr#1\crcr}}\right.}

% 插图、表题、脚注、横线
\long\def\topinsert{\par\pagebreak[1]\vskip 1.6em
  \noindent\begin{minipage}{\linewidth}\centering}
\long\def\endinsert{\par\end{minipage}\par\vskip 1.6em}
\newcount\figureno \figureno=0
\long\def\figurecaption#1{\global\advance\figureno by \@ne
  \par\vskip 5pt\begingroup\small
  \leftskip=2em\rightskip=2em
  \noindent 图~\the\figureno.\ #1\par\endgroup}
\newcount\tableno \tableno=0
\long\def\tablecaption#1{\global\advance\tableno by \@ne
  \par\vskip 3pt\begingroup\small\centering
  表~\the\tableno.\ #1\par\endgroup}
\def\thicktablerule{\rule{\linewidth}{0.8pt}}
\def\thintablerule{\rule{\linewidth}{0.4pt}}
\newcommand\sfn[2]{%
  \begingroup\renewcommand\thefootnote{#1}\footnotetext{#2}\endgroup
  \textsuperscript{#1}}
\makeatother

% ------------------------------------------------------------------
% 原文记号宏（luescher-macros.tex，同英文版）
% ------------------------------------------------------------------
\def\Bbb#1{\mathbb{#1}}
\def\frak#1{\mathfrak{#1}}
\renewcommand\Im{{\mathrm{Im}}\,}
\renewcommand\Re{{\mathrm{Re}}\,}
\let\ring\mathring
\newcommand\mib{\boldsymbol}
\providecommand\sixrm{\fontsize{6}{6.6}\selectfont}

% Definitions and abbreviations

% Roman letters in math formulae

\def\rmd{{\rm d}}
\def\rmD{{\rm D}}
\def\rme{{\rm e}}
\def\rmO{{\rm O}}
\def\rmU{{\rm U}}
\def\rmo{{\rm o}}

% Real and integer numbers

\def\rz{{\Bbb R}}
\def\gz{{\Bbb Z}}
\def\nz{{\Bbb N}}
\def\cz{{\Bbb C}}

% Special relations and symbols

\def\defeq{\mathrel{\mathop=^{\rm def}}}
\def\proof{\noindent{\sl Proof:}\kern0.6em}
\def\endproof{\hskip0.6em plus0.1em minus0.1em
\setbox0=\null\ht0=5.4pt\dp0=1pt\wd0=5.3pt
\vbox{\hrule height0.8pt
\hbox{\vrule width0.8pt\box0\vrule width0.8pt}
\hrule height0.8pt}}
\def\dual{\mathstrut^*\kern-0.1em}
\def\mod{\;\hbox{\rm mod}\;}
\def\lvec#1{\setbox0=\hbox{$#1$}
    \setbox1=\hbox{$\scriptstyle\leftarrow$}
    #1\kern-\wd0\smash{
    \raise\ht0\hbox{$\raise1pt\hbox{$\scriptstyle\leftarrow$}$}}
    \kern-\wd1\kern\wd0}
\def\rvec#1{\setbox0=\hbox{$#1$}
    \setbox1=\hbox{$\scriptstyle\rightarrow$}
    #1\kern-\wd0\smash{
    \raise\ht0\hbox{$\raise1pt\hbox{$\scriptstyle\rightarrow$}$}}
    \kern-\wd1\kern\wd0}

\def\slash#1{\setbox2=\hbox{$\displaystyle#1$}%
             \setbox3=\hbox{$\displaystyle/$}%
             #1\kern-0.8\wd2/\kern-1.0\wd3\kern0.8\wd2\kern0.5pt}
\def\ttilde#1{\setbox4=\hbox{$\displaystyle\tilde{\phantom{#1}}$}%
                    \hbox{$\displaystyle\tilde{\phantom{#1}}$}\kern-1.0\wd4%
         \raise1.6pt\hbox{$\displaystyle\tilde{\phantom{#1}}$}\kern-1.0\wd4%
         #1}
\def\sttilde#1{\setbox4=\hbox{$\scriptstyle\tilde{\phantom{#1}}$}%
                    \hbox{$\scriptstyle\tilde{\phantom{#1}}$}\kern-1.0\wd4%
         \raise1.0pt\hbox{$\scriptstyle\tilde{\phantom{#1}}$}\kern-1.0\wd4%
         #1}

\def\wick#1{\setbox2=\hbox{$\displaystyle#1$}
    \setbox3=\null\ht3=3.0pt\dp3=0.0pt\wd3=20.0pt
    #1\kern-\wd2\kern3.0pt\raise11.0pt\vbox{\hrule height0.3pt
    \hbox{\vrule width0.3pt\box3\vrule width0.3pt}}\kern-24.0pt\kern\wd2}

\def\longwick#1{\setbox2=\hbox{$\displaystyle#1$}
    \setbox3=\null\ht3=3.0pt\dp3=0.0pt\wd3=27.0pt
    #1\kern-\wd2\kern3.0pt\raise11.0pt\vbox{\hrule height0.3pt
    \hbox{\vrule width0.3pt\box3\vrule width0.3pt}}\kern-31.0pt\kern\wd2}

\def\verylongwick#1{\setbox2=\hbox{$\displaystyle#1$}
    \setbox3=\null\ht3=3.0pt\dp3=0.0pt\wd3=43.0pt
    #1\kern-\wd2\kern3.0pt\raise11.0pt\vbox{\hrule height0.3pt
    \hbox{\vrule width0.3pt\box3\vrule width0.3pt}}\kern-47.0pt\kern\wd2}


% Lattice derivatives

\def\nab#1{{\nabla_{#1}}}
\def\nabstar#1{{\nabla\kern0.5pt\smash{\raise 4.5pt\hbox{$\ast$}}
               \kern-5.5pt_{#1}}}
\def\drv#1{{\partial_{#1}}}
\def\drvstar#1{{\partial\kern0.5pt\smash{\raise 4.5pt\hbox{$\ast$}}
               \kern-6.0pt_{#1}}}
\def\ldrv#1{{\lvec{\,\partial}_{#1}}}
\def\ldrvstar#1{{\lvec{\,\partial}\kern-0.5pt\smash{\raise 4.5pt\hbox{$\ast$}}
               \kern-5.0pt_{#1}}}

% Units

\def\MeV{{\rm MeV}}
\def\GeV{{\rm GeV}}
\def\TeV{{\rm TeV}}
\def\fm{{\rm fm}}
\def\MSbar{\overline{\rm MS\kern-0.5pt}\kern0.5pt}

% Constants

\def\euler{\gamma_{\rm E}}

% Fields

\def\psibar{\overline{\psi}}
\def\chibar{\overline{\chi}}
\def\phibar{\overline{\phi}}
\def\lambdabar{\overline{\lambda}}
\def\ubar{\bar{u}}
\def\cbar{\bar{c}}
\def\dbar{\bar{d}}
\def\sbar{\bar{s}}
\def\Aax{A}
\def\Pax{P}
\def\Ptax{\tilde{\Pax}}

\def\ren#1{#1_{\hbox{\sixrm R}}}
\def\rgi#1{#1_{\hbox{\sixrm RGI}}}
\def\imp#1{#1_{\hbox{\sixrm I}}}
\def\reg#1{#1_{\rm reg}}
\def\lat#1{#1^{\rm lat}}
\def\rens#1#2{#1_{\hbox{\sixrm R},#2}}
\def\imps#1#2{#1_{\hbox{\sixrm I},#2}}

% Dirac matrices

\def\dirac#1{\gamma_{#1}}
\def\diracstar#1#2{
    \setbox0=\hbox{$\gamma$}\setbox1=\hbox{$\gamma_{#1}$}
    \gamma_{#1}\kern-\wd1\kern\wd0
    \smash{\raise4.5pt\hbox{$\scriptstyle#2$}}}
\def\dirachat{\hat{\gamma}_5}

% Gauge group

\def\SUtwo{{\rm SU(2)}}
\def\SUthree{{\rm SU(3)}}
\def\SUn{{\rm SU}(N)}
\def\sutwo{\frak{su}(2)}
\def\suthree{\frak{su}(3)}
\def\sun{\frak{su}(N)}
\def\tr{{\rm tr}}
\def\Tr{{\rm Tr}}
\def\Ad{{\rm Ad}\,}
\def\Psu{{\cal P}}
\def\Dsu{\frak{N}}

% Action

\def\SG{S_{\rm G}}
\def\SF{S_{\rm F}}
\def\Sw{S_{\rm w}}
\def\SGfl{S_{\rm G,fl}}
\def\SFfl{S_{\rm F,fl}}
\def\SFPfl{S_{\rm FP,fl}}
\def\Sms{S_{\rm ms}}
\def\Dw{D_{\rm w}}
\def\Dwdag{{\Dw}\kern-4pt^{\dagger}\kern1pt}
\def\csw{c_{\rm sw}}
%\def\cG{c_{\hbox{\sixrm G}}}
%\def\cF{c_{\hbox{\sixrm F}}}
\def\ZA{Z_A}
\def\ZP{Z_P}
\def\ZAh#1{\hat{Z}_A^{#1}}
\def\ZPh#1{\hat{Z}_P^{#1}}

% Dirac operator

\def\Dm{D}
\def\Dmdag{\Dm^{\dagger}\kern-1pt}
\def\Dms{D_s}
\def\Dmsdag{D_s\kern-2pt\vphantom{D}^{\dagger}\kern-1pt}

% Masses

\def\mpi{m_{\pi}}
\def\mK{m_K}
\def\mR#1{m_{\hbox{\sixrm R},#1}}

% Improvement

\def\csw{c_{\rm sw}}
\def\cfl{c_{\rm fl}}
\def\ca{c_A}
\def\cta{\tilde{c}_A}
\def\ctp{\tilde{c}_P}
\def\mq#1{m_{{\rm q},#1}}
\def\Mq{M_{\rm q}}

% Two-point correlation functions

\def\CPP#1{{\cal C}_{P}^{#1}}
\def\CPtP#1{{\cal C}_{\tilde{P}}^{#1}}
\def\CPhP#1{{\cal C}_{\hat{P}}^{#1}}
\def\CAP#1{{\cal C}_{\partial\kern-0.5pt A}^{#1}}
\def\CAtP#1{{\cal C}_{\partial\kern-0.5pt\tilde{A}}^{#1}}
\def\CdPP#1{{\cal C}_{\partial\partial\kern-0.5pt P}^{#1}}
\def\CQP#1{{\cal C}_{Q}^{#1}}

% Misc

\def\Nf{N_{\rm f}}
\def\eps{\epsilon}
\def\Mpi{M_{\pi}}
\def\Fpi{F_{\pi}}
\def\Gpi{G_{\pi}}
\def\Gpit{G_{\pi,t}}
\def\mK{m_{K}}
\def\FK{F_{K}}
\def\GK{G_{K}}
\def\GKt{G_{K,t}}
\def\CF{C_{\rm F}}
\def\Sigt{\Sigma_t}
\def\Sigtr{\Sigma_{t,\hbox{\sixrm R}}}
\def\gbar{\bar{g}}
\def\mbar{\kern0.5pt\overline{\kern-0.5pt m\kern-0.8pt}\kern0.8pt}
\def\zbar{\kern0.4pt\overline{\kern-0.4pt z\kern-0.6pt}\kern0.6pt}
\def\Zbar{\kern1.4pt\overline{\kern-1.4pt Z\kern-0.8pt}\kern0.8pt}
\def\Sigbar{\overline{\Sigma\kern-0.8pt}{\vphantom{\Sigma}}\kern0.8pt_t}
\def\cvec#1{\kern-0.5pt\vec{\kern0.5pt #1}}



\begin{document}

\rightline{\normalsize CERN-PH-TH/2013-025}
\vspace{1.2cm}
\begin{center}
{\LARGE\bfseries 手征对称性与杨--米尔斯梯度流}\\[0.8em]
{\normalsize Chiral symmetry and the Yang--Mills gradient flow}\\[1.2em]
{\large Martin L\"uscher}\\[0.8em]
{\itshape CERN, Physics Department, 1211 Geneva 23, Switzerland}\\[0.9em]
{\footnotesize JHEP 04 (2013) 123 \quad DOI: 10.1007/JHEP04(2013)123 \quad arXiv:1302.5246 [hep-lat]}
\end{center}
\vspace{0.8em}
\noindent\rule{\linewidth}{0.4pt}
\vspace{0.6em}

\noindent{\bfseries 摘要}\quad
近几年来，杨--米尔斯梯度流已被证明是非阿贝尔规范理论非微扰研究的一种有吸引力的工具。
本文考虑将该流向 QCD 夸克场的一个简单推广。
与纯规范梯度流的情形一样，
涉及正流时间局域场的关联函数的可重正化性，
可以借助一个 $4+1$ 维局域场论的表示来建立。
推广后的流在格点 QCD 中的应用包括：
非微扰重正化与 O($a$) 改善，
以及在手征极限下手征凝聚和赝标衰变常数的精确计算。

\vspace{0.6em}
\noindent\rule{\linewidth}{0.4pt}
\vspace{1.2em}


\lussec 1. 引言

鉴于其重正化性质
[\cite{WilsonFlow},\cite{RenFlow}]，
并且它在格点规范理论中的应用在技术上十分直接，
杨--米尔斯梯度流使我们能以许多有趣的方式探究非阿贝尔
规范理论的动力学。
例如，该流可用于精确的标度设定，
而且它给出了拓扑（瞬子）扇区如何在
格点 QCD 的连续极限中显现的确切理解
[\cite{WilsonFlow}]\kern1pt\sfn{$\dagger$}{%
在格点规范理论中，当流方程右端与 Wilson 方格作用量的梯度一致时，
该流常被称为 ``Wilson 流''。
本文在连续理论与格点上均使用``梯度流''一词。}。
此外，正流时间下的观测量
是非微扰重正化与步长标度研究中
自然而然要考虑的量
[\cite{StepScaling}--\cite{FritzschRamos}]。

流方程中可以包含物质场，也可以不包含。
本文考虑将流向 QCD 夸克场的一个相当平凡的推广：
规范场的流方程保持不变，
而夸克场随流时间的演化
由一个规范协变的热方程决定（见第 2 节）。
夸克场的引入并不会使流的
理论分析与实际实现显著复杂化。特别地，
仿照文献~[\cite{RenFlow}]，可以证明：一旦 QCD 的参数按通常方式重正化，
正流时间下规范不变的局域场关联函数的重正化
只需要对随时间演化的夸克场作
一个乘法重正化。

作为示例，本文具体完成了推广流的两个应用：
其一是格点 QCD 中轴矢流重正化常数计算的
一种新策略，
其二是在很大程度上通过计算正流时间下
标量夸克密度的期望值来得到手征凝聚。
在这两个例子中，流的使用在技术上都被证明是有吸引力的。
例如，手征凝聚可以很容易地以高精度获得，
因为不需要加性重正化。

在接下来的两节中，我们在连续理论中讨论梯度流
向夸克场的推广。
由于流方程保持手征对称性，
随时间演化场的关联函数
满足简单的手征 Ward 恒等式。
如第 4 节所解释的，这些恒等式使人们能把
这些关联函数与轻赝标介子的物理联系起来。
随后是论文更偏技术的部分（第 5--7 节），
其中在格点 QCD 的框架下建立该流。特别地，
我们将证明与之相联系的 $4+1$ 维场论
容许一个良定义的局域格点正规化。
最后，通过在 $2+1$ 味 QCD 中的一个示例计算
（第 8、9 节）演示所提出的流的应用
在数值格点 QCD 中的可行性。


\lussec 2. QCD 中的梯度流

本文所考虑的理论是规范群为 $\SUn$、含两味或更多有质量夸克
（处于规范群基础表示的多重态）的 QCD。
不过许多结果具有相当的普遍性，并不局限于 QCD。
理论在欧几里得空间中建立，
并照例通过泛函积分量子化。
在本节及下一节中采用维数正规化。
记号约定总结于附录 A。


\lussubsec 2.1 流方程

设 $A_{\mu}(x)$ 为 $\SUn$ 规范势，
$\psi(x)$ 为夸克场，$\psibar(x)$ 为在泛函积分中被积分掉的
反夸克场。后者带有狄拉克、颜色与味道指标，
为简单起见通常将其省略。

杨--米尔斯梯度流把规范场作为一个参数 $t\geq0$ 的函数来演化，
该参数被称为流时间。
从基本规范场出发，
\equation{
  \left.B_{\mu}\right|_{t=0}=A_{\mu},
  \enum
}
随时间演化的场 $B_{\mu}(t,x)$ 由微分方程
\equation{
  \partial_tB_{\mu}=D_{\nu}G_{\nu\mu},
  \enum
  \nexteq{2ex}
  G_{\mu\nu}=\partial_{\mu}B_{\nu}-\partial_{\nu}B_{\mu}+[B_{\mu},B_{\nu}],
  \qquad
  D_{\mu}=\partial_{\mu}+[B_{\mu},\,\cdot\;]
  \enum
}
确定（关于该专题的导引见文献~[\cite{WilsonFlow}]）。

正如第 1 节已经提到的，本文所考虑的流向夸克场的
推广是一个极小化的推广：
规范场的演化保持不变。
然而规范场出现在如下流方程之中\kern1pt\sfn{$\dagger$}{%
任何微分算符 $\Delta$ 的左作用
（在格点场论中则是差分算符的左作用）
通过要求关系式
$\eta^{\dagger}\lvec{\Delta}=(\Delta\eta)^{\dagger}$
对所有复值场 $\eta$ 成立来定义。}
\equation{
  \partial_t\chi=\Delta\chi,
  \qquad
  \partial_t\chibar=\chibar\lvec{\Delta},
  \enum
  \nexteq{2.5ex}
  \Delta=D_{\mu}D_{\mu},\qquad D_{\mu}=\partial_{\mu}+B_{\mu},
  \enum
}
它与初始条件
\equation{
  \left.\chi\right|_{t=0}=\psi,
  \qquad
  \left.\chibar\right|_{t=0}=\psibar,
  \enum
}
一起定义了随时间演化的夸克场与反夸克场
$\chi(t,x)$ 与 $\chibar(t,x)$。
在流方程 (2.4) 中，
规范协变的拉普拉斯算子 $\Delta$
例如可以换成狄拉克算子 $\slash{D}$ 的平方，
或乘以一个比例因子，
但这类修改似乎并无任何好处。

本节引入的夸克场流
与很多年前文献~[\cite{Guesken},\cite{AlexandrouEtAl}]
提出的源光滑化操作类似。
然而相对于这些流行的方法，
存在若干重要差别，
其一是该流在四维时空中运作，
而不是作用在等时超平面上的场上。
此外，流时间是连续变化的，
且规范场并不被固定为基本规范场，
而是与夸克场一同演化。


\lussubsec 2.2 局域复合场

由于流方程是规范协变的，
随时间演化的场在规范变换下与基本场按同样方式变换。
既在时空局域又在流时间局域的规范不变复合场的例子
是如下密度：
\equation{
  E_t(x)=-\frac{1}{2}\tr\{G_{\mu\nu}(t,x)G_{\mu\nu}(t,x)\},
  \enum
  \nexteq{2.5ex}
  S^{rs}_t(x)=\chibar_r(t,x)\chi_s(t,x),
  \qquad
  P^{rs}_t(x)=\chibar_r(t,x)\dirac{5}\chi_s(t,x),
  \enum
}
其中 $r,s$ 为味道指标。

通过初始条件 (2.1)、(2.6)，
随时间演化的场依赖于基本场。
因此从 QCD 泛函积分的观点看，
像密度 (2.7)、(2.8) 这样的复合场是可观测量，
类似于 Wilson 圈或普通的赝标密度等。
下文中，我们关心的是这些场的关联函数，
即由任意流时间的基本场构成的局域场乘积的期望值。


\lussubsec 2.3 微扰论

如文献~[\cite{WilsonFlow},\cite{RenFlow}] 所解释的，
这类关联函数可以直接按规范耦合的幂次展开。
首先将流方程按基本场的幂次求解。
把这些展开代入局域场，
便得到基本场关联函数的线性组合，
后者可用标准 QCD 费曼规则计算。

在微扰论中，
为了避免一些技术上的微妙之处，
将流方程 (2.2) 与 (2.4) 换成
\equation{
  \partial_tB_{\mu}=D_{\nu}G_{\nu\mu}+\alpha_0D_{\mu}\partial_{\nu}B_{\nu},
  \enum
  \nexteq{2.5ex}
  \partial_t\chi=\Delta\chi-
  \alpha_0\partial_{\nu}B_{\nu}\chi,
  \qquad
  \partial_t\chibar=\chibar\lvec{\Delta}+
  \alpha_0\chibar\partial_{\nu}B_{\nu},
  \enum
}
与参数 $\alpha_0>0$ 成比例的各项起到阻尼规范模式的作用，
并且可以通过一个依赖时间的规范变换将其去掉
[\cite{WilsonFlow}]。
因此规范不变场的关联函数不受这些额外项的影响。

将夸克流方程按其积分形式迭代，
便得到随时间演化的夸克场 $\chi(t,x)$
按基本场幂次的展开：
\equation{
  \chi(t,x)=\int\rmd^{D}\kern-1pt y\left\{K_t(x-y)\psi(y)
  +\int_0^t\rmd s\,K_{t-s}(x-y)\Delta'\chi(s,y)\right\},
  \enum
  \nexteq{3.0ex}
  \Delta'=(1-\alpha_0)\partial_{\nu}B_{\nu}
  +2B_{\nu}\partial_{\nu}+B_{\nu}B_{\nu},
  \enum
}
随时间演化的规范场也有相应的积分方程
[\cite{WilsonFlow}]，其中
\equation{
  K_t(z)={\rme^{-z^2/4t}\over(4\pi t)^{D/2}}
  \enum
}
表示 $D$ 维拉普拉斯算子的热核。
在规范耦合的最低阶，
可略去式 (2.11) 中的相互作用项，
于是得到随时间演化夸克场两点函数的表达式
\equation{
  \langle\chi(t,x)\chibar(s,y)\rangle=
  \int{\rmd^{D}\kern-1pt p\over(2\pi)^D}\,
  \rme^{ip(x-y)}
  {\rme^{-(t+s)p^2}\over M_0+i\slash{p}}+\rmO(g_0^2)
  \enum
}
其中 $g_0$ 与 $M_0$ 分别为裸耦合与（对角的）夸克质量矩阵。
由这些方程可以明显看出夸克流的光滑化特征。
此外可见，
领头阶下的光滑化范围 $\sqrt{8t}$ 与规范场情形相同。

在单圈微扰论中，
有 $8$ 个自能图对两点函数 (2.14) 有贡献。
这些图中的顶点来自 QCD 作用量
以及积分方程 (2.11) 的迭代。
按照文献~[\cite{RenFlow}] 的思路，
容易算出各图的紫外发散部分，
并且随即发现：只需照常重正化夸克质量，
并按下式重正化场，
两点函数即可被重正化：
\equation{
  \chi=Z_{\chi}^{-1/2}\ren{\chi},
  \qquad
  \chibar=\ren{\chibar}Z_{\chi}^{-1/2}.
  \enum
}
在 $D=4-2\eps$ 维的 $\MSbar$ 方案中，
对该场重正化常数的计算给出
\equation{
  Z_{\chi}=1+{3\CF\over16\pi^2\eps}g^2+\rmO(g^4),
  \qquad
  \CF={N^2-1\over2N},
  \enum
}
其中 $g$ 为重正化耦合。


\lussubsec 2.4 非零流时间下的夸克凝聚

在非零流时间下，
积分 (2.14) 中的大动量被指数压低，
因而当 $(t,x)\to(s,y)$ 时夸克两点函数没有奇异性。
短距离奇异性的缺失是正流时间下关联函数的一般特征，
它源于流方程的光滑化性质。
特别地，``随时间演化的夸克凝聚''
\equation{
  \Sigt^{rr}=-\langle S^{rr}_t(x)\rangle
  \enum
}
不需要加性重正化。

在微扰论中，
重正化的凝聚
\equation{
  \rens{\Sigma}{t}^{rr}=Z_{\chi}\Sigt^{rr}
  \enum
}
可以按重正化耦合 $g$ 的幂次算出，
其系数依赖于重正化夸克质量 $\mR{r},\mR{s},\ldots$
与流时间 $t$。
四维下的领头阶项由下式给出：
\equation{
  \left.\rens{\Sigma}{t}^{rr}\right|_{g=0}=
  {2N\mR{r}\over(4\pi)^2t}
  \int_0^{\infty}\rmd v\,{\rme^{-vz}\over(1+v)^2},
  \qquad z=2t\mR{r}^2,
  \enum
}
因而在手征极限下为零，
这与手征对称性只能在非微扰层面发生自发破缺这一事实相一致。


\lussec 3. $\mib D$+1 维场论

以下各节的讨论在很大程度上基于这样一个事实：
随时间演化场的关联函数与
$D$+1 维局域场论中的关联函数一致，
多出的一维即流时间。
这一观察可追溯到 Zinn--Justin 与 Zwanziger 关于
朗之万方程的开创性工作 [\cite{Zinn},\cite{ZinnZwanziger}]，
并且最近使得（纯规范的）梯度流的重正化性
能够在微扰论的所有阶得以建立
[\cite{RenFlow}]。

如下文将变得清楚的，
在 $D$+1 维理论中纳入夸克场是直截了当的。
总体框架与文献~[\cite{RenFlow}] 完全相同，
那里可以找到对该专题的导引及更多细节。


\lussubsec 3.1 场与作用量

除已经遇到的基本场和随时间演化的场之外，
$D$+1 维理论还涉及拉格朗日乘子场
$L_{\mu}(t,x)$、$\lambda(t,x)$ 与 $\lambdabar(t,x)$。
后者是与夸克场带有相同指标的费米子场，
而 $iL_{\mu}(t,x)$ 是取值于 $\SUn$ 李代数的矢量场。

如文献~[\cite{RenFlow}] 所解释的，
微扰论中的规范固定要求引入法捷耶夫--波波夫鬼场
$c(x),\cbar(x)$ 及相应的随时间演化场 $d(t,x),\dbar(t,x)$，其中
\equation{
  \left.d\right|_{t=0}=c.
  \enum
}
场 $\dbar(t,x)$ 不施加边界条件；
它在许多方面起着类似于拉格朗日乘子场的
作用。
除边界条件 (2.1)、(2.6) 与 (3.1) 外，
此处所有场均被假定为不受约束的。

理论的 total 作用量
\equation{
  S_{\rm tot}=S+\SGfl+\SFPfl+\SFfl,
  \enum
}
包含流时间为零的场的规范固定 QCD 作用量 $S$
以及体作用量
\equation{
  \SGfl=-2\int_0^{\infty}\rmd t\int\rmd^D\kern-1pt x\, 
  \tr\bigl\{L_{\mu}(t,x)
  \bigl(\partial_tB_{\mu}-
        D_{\nu}G_{\nu\mu}-\alpha_0D_{\mu}\partial_{\nu}B_{\nu}
  \bigr)(t,x)\bigr\},
  \enum
  \nexteq{3.0ex}
  \SFPfl=-2\int_{0}^{\infty}\rmd t\int\rmd^D\kern-1pt x\,
  \tr\bigl\{
  \dbar(t,x)\bigl(\partial_td-\alpha_0 D_{\mu}\partial_{\mu}d\bigr)(t,x)
  \bigr\},
  \enum
  \nexteq{3.0ex}    
  \SFfl=\int_0^{\infty}\rmd t\int\rmd^D\kern-1pt x\,
  \bigl\{\lambdabar(t,x)
         (\partial_t-\Delta+\alpha_0\partial_{\nu}B_{\nu})\chi(t,x)
  \noenum
  \nexteq{1.5ex}
  {\phantom{\SFfl=\int_0^{\infty}\rmd t\int\rmd^D\kern-1pt x\,\bigl\{}}
  +\chibar(t,x)
  (\lvec{\partial}_t-\lvec{\Delta}-\alpha_0\partial_{\nu}B_{\nu})
  \lambda(t,x)\bigr\}.
  \enum
}
相对于文献~[\cite{RenFlow}] 中研究的纯规范流，
体夸克作用量 $\SFfl$ 是 total 作用量中唯一的新增项。


\lussubsec 3.2 关联函数

规范场、夸克场、鬼场与拉格朗日乘子场的
$n$ 点关联函数，
形式上通过带``玻尔兹曼因子''
$\exp\{-S_{\rm tot}\}$ 的对全体场的泛函积分来定义。

对于规范场与夸克场的关联函数，
泛函积分可以在一定程度上通过先对拉格朗日乘子场
积分而算出。
由于作用量对这些场是线性的，
积分给出泛函狄拉克 $\delta$ 函数的乘积。
体作用量的选取使得这些 $\delta$ 函数
等价于把流方程强加于随时间演化的规范场与夸克场。
除数学上的微妙之处以及可能非平凡的雅可比因子外，
这样得到的关联函数与第 2 节中直接计算的结果一致：
即用初始条件 (2.1)、(2.6) 求解流方程，
并把算出的随时间演化场当作 QCD 可观测量。

$D$+1 维理论与直接计算关联函数之间的等价性，
可在微扰论中 [\cite{RenFlow}]
及格点 QCD 中（第 5 节）严格证明。
特别地，所有不涉及随时间演化夸克场
或相应拉格朗日乘子场的关联函数，
都与文献~[\cite{RenFlow}] 中研究的（夸克场不随流时间演化的）
关联函数一致。
$D$+1 维理论的这一``约化性质''
强烈限制了重正化所需的可能抵消项的形式
（见 3.4 小节）。


\lussubsec 3.3 费米子积分

由于 $D$+1 维中的作用量对费米子场是二次的，
这些场的泛函积分可以通过应用 Wick 定理求值（固定玻色场）。
基本 Wick 收缩值得明确写出，
因为这些表达式提供了对 $D$+1 维理论结构的洞察，
而且无论如何以后会经常用到。

基本场的 Wick 收缩由作用量和边界条件 (2.6) 决定。
特别地，流时间为零处夸克场的传播子
\equation{
  \wick{\psi(x)\psibar(y)}=S(x,y),
  \enum
  \nexteq{2.5ex}
  (\slash{D}+M_0)S(x,y)=\delta(x-y),
  \enum
}
与 $D$ 维中的传播子一致。
除夸克传播子 $S(x,y)$ 外，
随时间演化费米子场的收缩
\equation{
  \longwick{\lambda(t,x)\lambdabar(s,y)}=0,
  \enum
  \nexteq{2.5ex}
  \longwick{\chi(t,x)\lambdabar(s,y)}=\theta(t-s)K(t,x;s,y),
  \enum
  \nexteq{2.5ex}
  \longwick{\lambda(t,x)\chibar(s,y)}=\theta(s-t)K(s,y;t,x)^{\dagger},
  \enum
  \nexteq{2.5ex}
  \longwick{\chi(t,x)\chibar(s,y)}=\int\rmd^Dv\,
  \rmd^Dw\,
  K(t,x;0,v)S(v,w)K(s,y;0,w)^{\dagger},
  \enum
}
涉及夸克流方程的基本解
\equation{
  \{\partial_t-\Delta+
  \alpha_0\partial_{\nu}B_{\nu}\}K(t,x;s,y)=0\quad\hbox{if}\quad t\geq s,
  \enum
  \nexteq{2.0ex}
  \lim_{t\to s}K(t,x;s,y)=\delta(x-y),
  \enum
}
注意核 $K(t,x;s,y)$ 与味道无关，
正比于单位狄拉克矩阵，
且是颜色空间中的 $N\times N$ 复矩阵
（式 (3.10)、(3.11) 中核的伴随仅在指标空间上取）。

与流时间为零处的边界条件一致，
涉及基本夸克场与反夸克场的收缩
\equation{
  \longwick{\chi(t,x)\psibar(y)}=\int\rmd^Dv\,K(t,x;0,v)
  S(v,y),
  \enum
  \nexteq{2.5ex}
  \wick{\psi(x)\chibar(s,y)}=\int\rmd^Dw\,S(x,w)
  K(s,y;0,w)^{\dagger},
  \enum
  \nexteq{2.5ex}
  \longwick{\lambda(t,x)\psibar(y)}=
  \wick{\psi(x)\lambdabar(s,y)}=0,
  \enum
}
作为特例已包含在式 (3.9)--(3.11) 中。
然而要注意，
时空边界附近关联函数的行为
往往会因短距离奇异性而复杂化
[\cite{Symanzik}]。
例如极限
\equation{
  \lim_{s\to0}\longwick{\lambda(0,x)\chibar(s,y)}=
  \lim_{t\to0}\longwick{\chi(t,x)\lambdabar(0,y)}=
  \delta(x-y),
  \enum
}
在 $t=s=0$ 处与收缩 (3.16) 相差一个接触项。


\lussubsec 3.4 重正化

在微扰论中，
$D$+1 维理论的重正化性质可以用幂计数
与对称性论证来确定。
由于已知只让规范场随流时间演化的理论，
其重正化不需要超出通常的 QCD 参数与场重正化
[\cite{RenFlow}]，
找出作用量全部可能的局域抵消项的工作因此得到简化。
特别地，在流时间 $t>0$ 处，
场 $B_{\mu}(t,x)$ 与 $L_{\mu}(t,x)$ 无需重正化。

当随时间演化的夸克场被纳入理论时，
可能需要向 total 作用量添加进一步的抵消项。
然而 3.2 小节末提到的约化性质意味着，
任何这样的项都必须正比于夸克场或
相应的拉格朗日乘子场。
此外如文献~[\cite{RenFlow}] 的 7.2 小节所解释的，
体抵消项（即在全部 $D$+1 个坐标上积分的局域项）是被排除的，
因为大流时间下的关联函数由树图给出。

考虑到未重正化理论的对称性，
以及夸克场与拉格朗日乘子场的幂计数量纲分别为 $3/2$ 和 $5/2$，
稍加思考即可发现唯一可能新增的抵消项正比于
\equation{
  \int\rmd^D\kern-1pt x\,
  \left\{\lambdabar(0,x)\psi(x)+\psibar(x)\lambda(0,x)\right\}.
  \enum
}
该项在费曼图中的插入，
等价于正流时间下费米子场
以及类似地反费米子场的一个重正化
\equation{
  \chi=Z_{\chi}^{-1/2}\ren{\chi},
  \qquad
  \lambda=Z_{\chi}^{1/2}\ren{\lambda},
  \enum
}
因此，除 QCD 参数与场的通常重正化之外，
$D$+1 维理论的重正化还要求按上述方式
重正化随时间演化的夸克场及相应的拉格朗日乘子场。


\lussec 4. 手征对称性关系

由于夸克流方程 (2.4) 保持手征对称性，
像密度 (2.8) 这样的随时间演化复合场
在手征转动下以简单的方式变换，
从而使我们能以有用的新方式研究该对称性的自发破缺。
本节的主要目标，是把随时间演化的凝聚 (2.17)
与手征极限下的普通手征凝聚
及赝标衰变常数联系起来。


\lussubsec 4.1 夸克场方程与 PCAC 关系

现在把时空维数取为 4，
并暂时忽略正规化与重正化的需要，
也就是说，严格地讲本小节的论述只在微扰论的树图级成立。

或许有些出乎意料，
$4+1$ 维理论中的夸克场方程
\equation{
  \langle\{(\slash{D}+M_0)\psi(x)-\lambda(0,x)\}\kern1pt
  \phi_1(t_1,x_1)\ldots\phi_n(t_n,x_n)
  \rangle=
  \hbox{contact terms},
  \enum
}
与 QCD 中的场方程相差一项，
该项与流时间非零处的夸克场耦合。
式 (4.1) 对任意局域场 $\phi_1,\ldots,\phi_n$ 成立，
而当 $4+1$ 维中所有点 $(t_1,x_1),\ldots,(t_n,x_n)$
都不同于 $(0,x)$ 时接触项为零。
场方程 (4.1) 的正确性由恒等式
\equation{
  (\slash{D}+M_0)\wick{\psi(x)\psibar(y)}
  =\longwick{\lambda(0,x)\psibar(y)}+\delta(x-y),
  \enum
  \nexteq{2.5ex}
  (\slash{D}+M_0)\wick{\psi(x)\chibar(s,y)}
  =\longwick{\lambda(0,x)\chibar(s,y)},
  \enum
  \nexteq{2.5ex}
  (\slash{D}+M_0)\wick{\psi(x)\lambdabar(s,y)}
  =\longwick{\lambda(0,x)\lambdabar(s,y)},
  \enum
}
得出，它们是 3.3 小节所引费米子场 Wick 收缩公式的直接推论。
此外，式 (4.2)--(4.4) 表明接触项只来自基本夸克场的收缩。

现在令
\equation{
  \Aax^{rs}_{\mu}(x)=\psibar_r(x)\dirac{\mu}\dirac{5}\psi_s(x),
  \qquad
  \Pax^{rs}(x)=\psibar_r(x)\dirac{5}\psi_s(x),
  \enum
}
为 QCD 的轴矢流与密度。
对于 $r\neq s$，夸克场方程随即蕴含 PCAC 关系
\equation{
  \langle
  \{\partial_{\mu}\Aax^{rs}_{\mu}(x)-(m_{0,r}+m_{0,s})\Pax^{rs}(x)
  +\Ptax^{rs}(x)\}\kern1pt
  \phi_1(t_1,x_1)\ldots\phi_n(t_n,x_n)\rangle
  \noenum
  \nexteq{2.0ex}
  \quad=\hbox{contact terms},
  \enum
}
其中 $m_{0,r},m_{0,s}$ 是味道标记为 $r,s$ 的夸克的裸质量，而
\equation{
  \Ptax^{rs}(x)=\lambdabar_r(0,x)\dirac{5}\psi_s(x)
  +\psibar_r(x)\dirac{5}\lambda_s(0,x).
  \enum
}
特别地，恒等式
\equation{
  \langle
  \{\partial_{\mu}\Aax^{rs}_{\mu}(x)-(m_{0,r}+m_{0,s})\Pax^{rs}(x)
  +\Ptax^{rs}(x)\}\kern1pt P^{sr}_t(y)\rangle=0
  \enum
}
对所有流时间 $t>0$ 以及包括 $x=y$ 在内的所有点 $x$ 成立。
此情形中没有接触项，因为对式 (4.8) 左端关联函数有贡献的
费米子传播子在 $4+1$ 维中全部从 $(0,x)$ 走到 $(t,y)$，
或从 $(t,y)$ 走到 $(0,x)$。
因而收缩 (4.2) 不出现。
鉴于流方程的光滑化性质，
事实上该关联函数的所有部分都是 $x$ 的正规函数。

Wick 收缩结构的另一个纯代数推论是关系式
\equation{
  \int\rmd^4x\,\langle\Ptax^{rs}(x)P^{sr}_t(y)\rangle
  =\langle S^{rr}_t(y)\rangle+\langle S^{ss}_t(y)\rangle,
  \enum
}
将它与 PCAC 关系 (4.8) 相结合，
便得到一个恒等式
\equation{
  (m_{0,r}+m_{0,s})\int\rmd^4x\,\langle\Pax^{rs}(x)P^{sr}_t(y)\rangle
  =-\Sigt^{rr}-\Sigt^{ss}
  \enum
}
它使人们能把随时间演化的夸克凝聚与赝标介子的物理联系起来
（见 4.3 小节）。在推导此式的过程中，
假定夸克质量与流时间均为正，
否则在大 $x$ 处没有边界项、在 $x=y$ 处没有不可积奇异性
就无法得到保证。


\lussubsec 4.2 重正化的 PCAC 关系

超出微扰论树图级后，
式 (4.8)--(4.10) 中的场和夸克质量需要重正化。
这些等式于是成为重正化关联函数之间的关系，
其确切形式依赖于所选的归一化约定。
这里有一些自然的选择，
以下的目的就是把它们明确下来。
讨论隐含地假设了理论的合理正规化，
例如维数正规化或保持大多数对称性的 Wilson 格点表述
[\cite{Wilson}]。与前一样，只考虑味非对角道。

非单态轴矢流与密度作乘法重正化：
\equation{
  \rens{\Aax}{\mu}^{rs}=Z_A\Aax_{\mu}^{rs},
  \qquad
  \ren{\Pax}^{rs}=Z_P\Pax^{rs},
  \enum
}
随时间演化的密度也是如此：
\equation{
  \rens{S}{t}^{rs}=Z_{\chi}S_t^{rs},
  \qquad
  \rens{P}{t}^{rs}=Z_{\chi}P_t^{rs}.
  \enum
}
如前所述，由于正流时间处不存在短距离奇异性，
这些场（包括味对角的那些）按其夸克内容重正化，
即如同 $4+1$ 维中非零距离处
随时间演化的夸克场与反夸克场的乘积那样。

对于场 $\Ptax^{rs}$，
其乘法可重正化性
\equation{
  \ren{\Ptax}^{rs}=\tilde{Z}_P\Ptax^{rs},
  \enum
}
不那么明显，
因为该场量纲为 4，可能与量纲 3 和 4 的其他场混合。
然而涉及拉格朗日乘子场 $\lambda,\lambdabar$
的 Wick 收缩具有这样的性质：
除接触项外，
$\Ptax^{rs}$ 与任何由流时间为零的规范场和夸克场构成的局域场的
两点关联函数为零。
因此被这类场所致的发散加性重正化被排除。
剩下只可能涉及拉格朗日乘子场的场，
但 $\Ptax^{rs}$ 是其中唯一具有量纲 4 及所需对称性质的场，
故加性重正化同样不可能，
于是该场必然作乘法重正化。

现在可以这样选取重正化场与重正化夸克质量 $\mR{r}$
的归一化，使得 PCAC 关系取如下形式：
\equation{
  \langle
  \{\partial_{\mu}\rens{\Aax}{\mu}^{rs}(x)-
  (\mR{r}+\mR{s})\ren{\Pax}^{rs}(x)
  +\ren{\Ptax}^{rs}(x)\}\kern1pt\rens{P}{t}^{sr}(y)\rangle=0.
  \enum
}
注意这一约定只固定了轴矢流 $\rens{\Aax}{\mu}^{rs}$、
乘积 $(\mR{r}+\mR{s})\ren{\Pax}^{rs}$
与密度 $\ren{\Ptax}^{rs}$ 三者之间的相对归一化。
不过后者的归一化可以通过要求方程
\equation{
  \int\rmd^4x\,\langle\ren{\Ptax}^{rs}(x)\rens{P}{t}^{sr}(y)\rangle
  =\langle\rens{S}{t}^{rr}(y)\rangle+\langle\rens{S}{t}^{ss}(y)\rangle
  \enum
}
以及恒等式
\equation{
  (\mR{r}+\mR{s})\int\rmd^4x\,
  \langle\ren{\Pax}^{rs}(x)\rens{P}{t}^{sr}(y)\rangle
  =-\rens{\Sigma}{t}^{rr}-\rens{\Sigma}{t}^{ss}
  \enum
}
（其中 $\rens{\Sigma}{t}^{rr}=Z_{\chi}\Sigma_t^{rr}$）成立而确定。
此后仍需对重正化夸克质量与随时间演化密度施加进一步的归一化条件，
但对质量可以采用某个标准约定，
而随时间演化场的归一化在我们关心的等式中往往相消，
故保持未定。

由式 (4.14)、(4.15) 隐含确定的轴矢流的归一化，
与通常选择的、使轴矢荷取整数值的归一化一致。
有不同办法可以说明这一点，
也许最直接的途径是从保持手征对称性的格点 QCD 表述出发
[\cite{GinspargWilson}--\cite{ExactChSy}]。
利用手征 Ward 恒等式 [\cite{BoulderReview}]，
可以证明在此情形下裸场满足式 (4.8)、(4.9)，
直至正比于格距而消失的项。
于是式 (4.14)、(4.15) 意味着 $Z_A=\tilde{Z}_P=1$，
从而给出重正化轴矢流的 canonical 归一化。
鉴于连续极限的普适性，
这一 canonical 归一化因此由这些等式保证，
而与所选的理论正规化无关。


\lussubsec 4.3 赝标矩阵元与手征凝聚

现在假定夸克多重态包含两个轻夸克，
称为 $u$、$d$ 夸克，二者具有相同的
重正化质量 $\ren{m}$。
若再无太多其余夸克，
$(u,d)$ 道的手征对称性预期会发生自发破缺，
于是存在一个赝标介子三重态——``$\pi$ 介子''，
其质量随 $\ren{m}\to0$ 而趋于零。

在大距离处，轻夸克赝标关联函数
由中间单 $\pi$ 介子态主导。
特别地，当 $x_0\to\infty$ 时
\equation{
  \int\rmd^3{\mib x}\,\langle\ren{\Pax}^{ud}(x)\ren{\Pax}^{du}(0)\rangle
  =-{\Gpi^2\over\Mpi}
  \rme^{-\Mpi x_0}\bigl\{1+\rmO(\rme^{-\Delta Ex_0})\bigr\},
  \enum
  \nexteq{2.5ex}
  \int\rmd^3{\mib x}\,\langle\rens{\Aax}{0}^{ud}(x)\ren{\Pax}^{du}(0)\rangle
  =\Fpi\Gpi\rme^{-\Mpi x_0}\bigl\{1+\rmO(\rme^{-\Delta Ex_0})\bigr\},
  \enum
}
其中 $\Delta E$ 表示到更高态的能量间隙，
而 $\Mpi$、$\Fpi$ 与 $\Gpi$
分别是 $\pi$ 介子质量、衰变常数
与轴矢密度的真空到 $\pi$ 介子矩阵元。

从四维理论的观点看，
$\rens{\Pax}{t}^{du}(x)$ 不是局域场，
但由于光滑核 $K(t,x;0,y)$ 在大距离 $|x-y|$ 处
近似按高斯型衰减，
它仍然是基本场的一个相当局域的表达式。
当 $x_0$ 远大于光滑范围 $\sqrt{8t}$ 时，
两点函数
\equation{
  \int\rmd^3{\mib x}\,\langle\ren{\Pax}^{ud}(x)\rens{\Pax}{t}^{du}(0)\rangle
  =-{\Gpi\Gpit\over\Mpi}
  \rme^{-\Mpi x_0}\bigl\{1+\rmO(\rme^{-\Delta Ex_0})\bigr\}
  \enum
}
因而与普通赝标关联函数按同样方式衰减。
系数 $\Gpit$ 等于某个算符的真空到 $\pi$ 介子矩阵元，
但此处把它视为由式 (4.19) 定义的量。

表征（两味）手征极限的低能有效常数
\equation{
  \Sigma=\lim_{m_{\rm R}\to0}\Fpi\Gpi,
  \enum
  \nexteq{2.5ex}
  F=\lim_{m_{\rm R}\to0}\Fpi,
  \enum
}
是手征凝聚 $\Sigma$
以及在 $\ren{m}=0$ 处的 $\pi$ 介子衰变常数值 $F$。
如前所述，
通过手征 Ward 恒等式 (4.16) 可以把它们与随时间演化的
轻夸克凝聚 $\rens{\Sigma}{t}=\rens{\Sigma}{t}^{uu}$
联系起来。利用 PCAC 关系
\equation{
  2\ren{m}\Gpi=\Mpi^2\Fpi,
  \enum
}
后者可以写成如下形式：
\equation{
  \rens{\Sigma}{t}=-{\Mpi^2\Fpi\over2\Gpi}
  \int\rmd^4x\,\langle \ren{P}^{ud}(x)\rens{P}{t}^{du}(0)\rangle.
  \enum
}
此式右端的积分在 $\ren{m}\to0$ 时发散，
其渐近行为由动量空间中关联函数的 $\pi$ 极点决定。
回忆式 (4.19)，即可得出结论
\equation{
  \Sigma=\lim_{m_{\rm R}\to0}{\rens{\Sigma}{t}\Gpi\over\Gpit},
  \enum
  \nexteq{2.5ex}
  F=\lim_{m_{\rm R}\to0}{\rens{\Sigma}{t}\over\Gpit}.
  \enum
}
于是在手征极限下，
随时间演化的凝聚 $\rens{\Sigma}{t}$
提供了手征对称性自发破缺的一个序参量，
它与标准凝聚 $\Sigma$ 只差一个比例常数。

式 (4.24)、(4.25) 的一个显著特点是：
其中完全没有提及轴矢流，
甚至没有通过归一化条件隐含涉及；
而且这些等式对任何（正的）流时间 $t$ 都成立。
此外，重正化常数 $Z_{\chi}$ 消失了，
因此例如衰变常数 $F$
可以直接由裸的随时间演化凝聚
与赝标真空到 $\pi$ 介子矩阵元给出。


\lussubsec 4.4 手征微扰论

在手征极限附近，
QCD 中许多量的渐近行为可以用手征微扰论描述。
如果 $\pi$ 介子的康普顿波长远大于梯度流的光滑范围，即若
\equation{
  8t\Mpi^2\ll1,
  \enum
}
则随时间演化的密度 $\rens{S}{t}^{rs}$、$\rens{P}{t}^{rs}$
从 $\pi$ 介子物理的观点看是局域场，
具有与流时间为零时的密度相同的对称性质。
因此在低能下，
随时间演化场的关联函数以及普通轴矢流与密度的关联函数
预期由 $\pi$ 介子中间态主导，
从而应当可以在手征微扰论框架内得到定量描述。

事实上，只要在手征拉氏量中加入适当的源项
[\cite{ChPT}]，
就可以很容易地把随时间演化的密度纳入手征微扰论。
在手征展开的每一阶，这些项要求引入若干有效耦合，
它们一般会依赖流时间。
于是可以算出手征凝聚 $\rens{\Sigma}{t}$
与矩阵元 $\Gpit$ 的展开，
从而洞察手征极限 (4.24)、(4.25)
究竟是如何达到的。


\lussec 5. 格点正规化

本节的出发点是 Wilson [\cite{Wilson}] 所提议的
把夸克场放到格点上的一种格点 QCD 表述。
对格点狄拉克算符或规范作用量不作特殊假设，
但二者应当是局域的，
并保持标准 Wilson 理论所保持的全部对称性。
为记号方便，取格距 $a=1$。
除非另有说明，格点在所有方向上被视为无限延伸。


\lussubsec 5.1 流方程

在格点上，随时间演化的规范场 $V(t,x,\mu)$ 由
\equation{
  \left.V\right|_{t=0}=U,
  \enum
  \nexteq{2.5ex}
  \partial_tVV^{-1}=-g_0^2\partial\Sw,
  \enum
}
定义，其中 $U(x,\mu)$ 是基本格点规范场，
$(\partial\Sw)(V,x,\mu)$ 是 Wilson 方格作用量的梯度
（进一步的解释见附录 A 及文献~[\cite{WilsonFlow}]）。

对于夸克场，边界条件与流方程的形式与连续理论相同，
只是式 (2.4) 中的算符 $\Delta$
换成格点拉普拉斯算子
\equation{
  \Delta=\nabstar{\mu}\nab{\mu},
  \enum
}
这里 $\nab{\mu}$ 与 $\nabstar{\mu}$
是在随时间演化规范场存在下的规范协变前向与后向差分算符。

如文献~[\cite{WilsonFlow}] 中已指出的，
流方程 (5.2) 解的整体存在性、唯一性与光滑性是被严格保证的。
由于 $\Delta$ 是有界算符，且它随流时间光滑变化，
可以证明同样的结论对夸克流方程也成立
[\cite{ArnoldI}]。
因此非零流时间的局域场——例如前面各节所考虑的那些场——
在格点上是完全良定义的。


\lussubsec 5.2 $4+1$ 维格点理论

与连续理论类似，
随时间演化局域场的格点 QCD 关联函数
可以从一个 $4+1$ 维格点场论得到。
为了能够完全控制这一构造，
先对该场论就流方程的一个离散化版本建立之，
其中流时间被限制在一个有限集合
\equation{
  0,\eps,2\eps,3\eps,\ldots,T
  \enum
}
上，各值由时间步长 $\eps>0$ 分隔。

离散演化方程
\equation{
  V(t+\eps,x,\mu)V(t,x,\mu)^{-1}=\rme^{-\eps g_0^2(\partial\Sw)(V,x,\mu)}
  \enum
}
对规范场的积分，相当于用欧拉方法
（或者等价地，通过施加一列 ``stout 涂抹'' 步骤 [\cite{Stout}]）
对连续方程 (5.2) 作近似积分。
对夸克场而言，
离散化的流方程形式与连续情形相同，
只是夸克场的时间导数换成前向差分
\equation{
  \partial_t\chi(t,x)={1\over\eps}\{\chi(t+\eps,x)-\chi(t,x)\}.
  \enum
}
注意这些离散方程可以逐步求解：
从 $t=0$ 的基本场出发，
递推地从时间 $t$ 推进到 $t+\eps$。
特别地，解由场的初值唯一确定。

在 $4+1$ 维格点理论中，
泛函积分中积分的所有场只在时刻 (5.4) 上定义。
由于此处不需要规范固定，
没有鬼场，作用量中也没有规范固定项，
但边界条件 (2.6)、(5.1) 照旧施加。
拉格朗日乘子场在时刻 $T$ 的分量与其他场退耦，
因此可置为零（或从一开始就略去）。

体作用量的一种可能选取是
\equation{
  \SGfl=-2\sum_{t=0}^{T-\eps}\sum_{x,\mu}
  \tr\bigl\{L(t,x,\mu)
  \noenum
  \nexteq{-0.5ex}
  {\phantom{\SGfl=-2\sum_{t=0}^{T-\eps}\sum_{x,\mu}}}
  \times
  \Psu\bigl(
  V(t+\eps,x,\mu)V(t,x,\mu)^{-1}-\rme^{-\eps g_0^2(\partial\Sw)(V,x,\mu)}
  \bigr)\bigr\},
  \enum
  \nexteq{2.5ex}
  \SFfl=\eps\sum_{t=0}^{T-\eps}\sum_x\bigl\{
  \lambdabar(t,x)\left(\partial_t-\Delta\right)\chi(t,x)
  +\chibar(t,x)\left(\lvec{\partial}_t-\lvec{\Delta}\right)\lambda(t,x)
  \bigr\},
  \enum
}
其中映射
\equation{
  M\to\Psu(M)={1\over2}(M-M^{\dagger})
  -{1\over2N}\kern1pt\tr(M-M^{\dagger})
  \enum
}
把任意复 $N\times N$ 矩阵映为 $\SUn$ 李代数的一个元素。
除这些作用量及格点 QCD 作用量外，
测度项
\equation{
  \Sms=\sum_{t=0}^{T-\eps}\sum_{x,\mu}
  {\cal L}_{\rm ms}(V(t+\eps,x,\mu)V(t,x,\mu)^{-1}),
  \enum
  \nexteq{3.0ex}
  {\cal L}_{\rm ms}(W)=
  \cases{-\ln\det(J(W)) & if $W\in\Dsu$,\cr
         \noalign{\vskip1ex}
         \infty         & otherwise,\cr}
  \enum
}
也必须包含在该格点理论的 total 作用量之中
（映射 (5.9) 的雅可比矩阵 $J(W)^{ab}$
及单位元的邻域 $\Dsu\subset\SUn$ 的定义见附录 B）。

给定了场及其作用量之后，
$4+1$ 维格点理论的泛函积分按标准方式定义。
注意测度项 (5.10)、(5.11)
实际上把泛函积分中的链变量限制在使
\equation{
  V(t+\eps,x,\mu)V(t,x,\mu)^{-1}\in\Dsu
  \quad\hbox{for all $t,x,\mu$}.
  \enum
}
成立的区域内。
由于作用量 $\SGfl$ 是纯虚的，
对拉格朗日乘子场 $L(t,x,\mu)$ 的积分不是绝对收敛的，
但应当理解，
由此可能产生的积分不确定性
通过在作用量中加入一个正比于拉格朗日乘子场平方的无穷小项来消除。


\lussubsec 5.3 四维理论的恢复

这样定义的、由场 $V,\chi,\chibar$ 构成的局域场的关联函数，
可以在一定程度上通过先对拉格朗日乘子场积分而求值。
由于 total 作用量对这些场是线性的，
积分给出一个乘积
\equation{
  \prod_{t=0}^{T-\eps}\prod_{x,\mu}
  \delta\bigl\{\Psu\bigl(
  V(t+\eps,x,\mu)V(t,x,\mu)^{-1}\bigr)-
  \Psu\bigl(\rme^{-\eps g_0^2(\partial\Sw)(V,x,\mu)}\bigr)\bigr\}
  \noenum
  \nexteq{1.0ex}
  \times\prod_{t=0}^{T-\eps}\prod_{x}
  \delta\bigl\{\partial_t\chi(t,x)-\Delta\chi(t,x)\bigr\}
  \delta\bigl\{\chibar(t,x)\lvec{\partial}_t-
  \chibar(t,x)\lvec{\Delta}\bigr\}
  \enum
}
即狄拉克 $\delta$ 函数的乘积。
当 $V,\chi,\chibar$ 满足离散流方程时，
这些 $\delta$ 函数的自变量为零；
而且对足够小的步长 $\eps$，
事实上没有别的办法能使这些自变量为零%
\kern1pt\sfn{$\dagger$}{%
若 $\exp\{\eps g_0^2(\partial\Sw)(V,x,\mu)\}\in\Dsu$
在格点的所有链上及所有流时间处成立，
则唯一性得到保证。
一些严格的估计表明：
若 $\eps\leq(24\sqrt{2}N)^{-1}$，则该条件被满足。}。
对规范场而言，
此处需要注意的要点是，
$\delta$ 函数中的映射 $\Psu$
作用在单位元小邻域内的 $\SUn$ 矩阵上。
附录 B 所引映射的性质于是意味着：
$\delta$ 函数所强加的约束等价于
把流方程 (5.5) 强加于规范场。

现在可以利用这些 $\delta$ 函数
对所有流时间 $t>0$ 处的场变量进行积分。
从时刻 $t=T$ 的场出发并向更小时间推进，
$\delta$ 函数允许递归地消去全部场，
直到剩下的积分变量只有基本场 $U,\psi,\psibar$。
在这一过程中，
规范场变量的 $\delta$ 函数积分产生映射 $\Psu$
的雅可比因子的乘积，
但这些因子被作用量中的测度项抵消
（见附录 B）。

计算于是表明：
不涉及拉格朗日乘子场的局域场的 $4+1$ 维关联函数，
与同一组场在四维中通过离散流方程和 QCD 泛函积分
定义的关联函数精确相等。


\lussubsec 5.4 连续时间极限

至此，关联函数中的随时间演化场仍是
由流方程的欧拉积分得到的那些场。
由于相关的积分误差为 $\eps$ 量级，
当 $\eps\to0$（且平凡地 $T\to\infty$）时
便恢复连续时间的关联函数。
此外，测度项
\equation{
  \Sms=\sum_{t=0}^{T-\eps}\sum_{x,\mu}
  \ln\det J\bigl(\rme^{-\eps g_0^2(\partial\Sw)(V,x,\mu)}\bigr)
  =\rmO(\eps)
  \enum
}
在此极限下为零。

下文中通常考虑格点理论的连续时间表述，
但 $4+1$ 维泛函积分总是理解为
通过离散时间泛函积分的极限 $\eps\to0$ 来定义。


\lussec 6. 壳上 O($\mib a$) 改善

涉及非零流时间处的场的关联函数对格距的依赖关系，
最好在 $4+1$ 维理论中讨论。
在这一框架下，
基于局域性、幂计数与对称性的标准论证可以照常应用。
相关的 Symanzik 局域有效理论
[\cite{SymanzikSeffI},\cite{SymanzikSeffII}]
与 $4+1$ 维连续理论一致，
只是在作用量和局域场上添加了 O($a$) 及更高阶的格距项。
此外，壳上改进 [\cite{OnShell}]
意味着格距效应的抵消只限于
$4+1$ 维中非零距离处的关联函数。

尽管理论上的讨论适用面更广，
从现在起假定采用格点 QCD 的 O($a$) 改进 Wilson 表述
[\cite{SW},\cite{SFimp}]。
在可能之处，沿用文献~[\cite{SFimp}] 的约定与记号。


\lussubsec 6.1 有效作用量

由于格点理论中的流方程在经典层面是 O($a$) 改进的，
又由于圈图对渐近大流时间下的关联函数没有贡献，
Symanzik 有效作用量中所有 O($a$) 项必须是边界项，
即流时间为零处的局域项。
再者，若有效理论只需描述壳上物理量——此处正是如此——
则许多项可以利用夸克场方程、流方程以及边界条件 (2.1)、(2.6)
消去。
在手征极限下，有效作用量中 O($a$) 项的一种可能选取是
\equation{
  a\int\rmd^4x\left\{c_1\psibar(x)\frac{i}{4}\sigma_{\mu\nu}F_{\mu\nu}(x)\psi(x)
  +c_2\lambdabar(0,x)\lambda(0,x)\right\},
  \enum
}
其中 $F_{\mu\nu}(x)$ 是基本规范势的场强张量
（此处及下文中，$c_1,c_2,\ldots$ 表示有效理论的系数）。

要描述非零夸克质量下的 O($a$) 格距效应，
所需的项要多得多。
在四维理论中，
这些项已被 Bhattacharya 等人 [\cite{ImpNonD}] 完全分类。
$4+1$ 维中的 Symanzik 有效作用量还额外包含一项
\equation{
  a\int\rmd^4x\left\{
  \lambdabar(0,x)(c_3M+c_4\kern0.5pt\tr M)\psi(x)+
  \psibar(x)(c_3M+c_4\kern0.5pt\tr M)\lambda(0,x)\right\}
  \enum
}
它对涉及体场的关联函数有贡献，
这里 $M$ 是有效理论中的夸克质量矩阵。
它对关联函数的影响等价于体费米子（及反费米子）场的一个重正化
\equation{
  \chi(t,x)\to Z\chi(t,x),
  \quad
  \lambda(t,x)\to Z^{-1}\lambda(t,x),
  \enum
  \nexteq{2.5ex}
  Z=1+ac_3M+ac_4\kern0.5pt\tr M,
  \enum
}

\lussubsec 6.2 有效局域场

在 Symanzik 理论中表示局域格点场的有效场的形式，
受到决定有效作用量形式的那些一般原理的限制。
特别地，表示正流时间处场的有效场没有 O($a$) 项。

不过这样的项在流时间为零的局域场情形中出现。
再次考虑只有无质量夸克的理论，
例如表示味道指标 $r\neq s$ 的轴矢流的有效场由下式给出：
\equation{
  (A_{\rm eff})_{\mu}^{rs}(x)
  =\psibar_r(x)\dirac{\mu}\dirac{5}\psi_s(x)+
   ac_5\partial_{\mu}\left\{\psibar_r(x)\dirac{5}\psi_s(x)\right\}
  \noenum
  \nexteq{2.0ex}
  {\phantom{(A_{\rm eff})_{\mu}^{rs}(x)=}}
  +ac_6\left\{\lambdabar_r(0,x)\dirac{\mu}\dirac{5}\psi_s(x)
  +\psibar_r(x)\dirac{\mu}\dirac{5}\lambda_s(0,x)
  \right\}+\ldots,
  \enum
}
相应的有效赝标密度
\equation{
  (P_{\rm eff})^{rs}(x)=
  \psibar_r(x)\dirac{5}\psi_s(x)
  \noenum
  \noenum
  \nexteq{2.0ex}
  {\phantom{(P_{\rm eff})^{rs}(x)=}}
  +
  ac_7\left\{\lambdabar_r(0,x)\dirac{5}\psi_s(x)
  +\psibar_r(x)\dirac{5}\lambda_s(0,x)
  \right\}+\ldots,
  \enum
}
也有类似的展开。

这些场的依赖质量的 O($a$) 修正
与四维理论中的相同 [\cite{ImpNonD}]。
但一般而言，
O($a$) 项的确定需要仔细考虑场之间可能发生的混合，
并可能导致相当复杂的表达式。


\lussubsec 6.3\/ {\rm O($a$)\kern-0.5pt} 改进的格点作用量

按照惯例，
依赖质量的 O($a$) 抵消项不在改进作用量中写出，
而是被纳入参数与场的重正化因子之中 [\cite{SFimp}]。
于是 $4+1$ 维中的 O($a$) 抵消项由下式给出：
\equation{
  \delta S_{\rm tot}=\sum_x\left\{\csw
  \psibar(x)\frac{i}{4}\sigma_{\mu\nu}\widehat{F}_{\mu\nu}(x)\psi(x)
  +\cfl\lambdabar(0,x)\lambda(0,x)\right\},
  \enum
}
其中 $\widehat{F}_{\mu\nu}(x)$ 表示格点上规范场强张量
$F_{\mu\nu}(x)$ 的标准三叶草表达式。
式 (6.7) 中的第一项（Sheikholeslami--Wohlert 项
[\cite{SW},\cite{SFimp}]）
是四维理论的改进本来就需要有的，
而正比于改善系数 $\cfl(g_0)$ 的那一项
则是为消去涉及体场的关联函数中的 O($a$) 格点效应所必需的。

给定了作用量之后，
O($a$) 改进理论中基本费米子场的 Wick 收缩
可以直接算出（见附录 C）。
在连续时间极限下，
这些收缩与连续理论中的收缩实质上相同，
只需作显然的修改
（时空坐标积分换成格点求和，
夸克传播子 $S(x,y)$ 换成有质量的 O($a$) 改进格点狄拉克算符之逆，
核 $K(t,x;s,y)$ 换成格点上夸克流方程的基本解）。
一个特例是随时间演化夸克场的收缩
\equation{
  \longwick{\chi(t,x)\chibar(s,y)}=\sum_{v,w}
  K(t,x;0,v)
  \left(S(v,w)-\cfl\delta_{vw}\right)
  K(s,y;0,w)^{\dagger}
  \enum 
}
它是唯一依赖改善系数 $\cfl$ 的收缩。
特别地，该收缩从而 $\psi$ 与 $\psibar$
场的关联函数都与 $\cfl$ 无关，
正如必须如此，
因为四维理论已通过纳入 Sheikholeslami--Wohlert 项而得到改进。

式 (6.8) 或许令人困惑的一点是极限
\equation{
  \lim_{t\to0}\kern0.5pt\longwick{\chi(t,x)\chibar(s,y)}=
  \wick{\psi(x)\chibar(s,y)}-\cfl K(s,y;0,x)^{\dagger}
  \enum
}
与 $\wick{\psi(x)\chibar(s,y)}$ 相差一个非接触项的
$a$ 阶项。
因此 O($a$) 改进所需的减除，
取决于关联函数中的夸克场处于零还是非零流时间。
这其实并不太令人惊讶，
因为这些场的重正化方式本就不同。
这两处差别只是反映了如下事实：
涉及随时间演化场的关联函数取连续极限时，
必须固定以物理单位给出的流时间；
而且只有按这种方式趋于连续极限时，
$4+1$ 维 Symanzik 有效理论才正确描述
格点理论对其连续极限的偏离。


\lussubsec 6.4 重正化的改进场

照例，加性重正化后的裸夸克质量 $\mq{r}$ 定义为
\equation{
  \mq{r}=m_{0,r}-m_{\rm c},
  \enum
}
其中 $m_{\rm c}$ 表示夸克质量简并理论中的临界裸质量。
引入相应的减除后裸质量矩阵 $\Mq$
以及组合
\equation{
   \mq{rs}=\frac{1}{2}\left(\mq{r}+\mq{s}\right)
   \enum
}
也是有帮助的。

正流时间处的格点费米子场需要乘法重正化，
以及通过一个依赖质量的因子所作的 O($a$) 改进
（参见 6.1 小节）。
夸克场与反夸克场用同一因子标度：
\equation{
  \rens{\chi}{r}=\{Z_{\chi}(1+b_{\chi}\mq{r}+
  \bar{b}_{\chi}\tr\Mq)\}^{1/2}\chi_r,
  \enum
}
而拉格朗日乘子场 $\lambda,\lambdabar$
用同一因子的倒数重正化%
\kern1pt\sfn{$\dagger$}{%
仿照文献~[\cite{SFimp},\cite{ImpNonD}]，
$b_X$ 与 $\bar{b}_X$ 泛指乘以依赖质量的 O($a$)
抵消项的改善系数。}。
于是像赝标密度这样的随时间演化复合场
\equation{
  \rens{P}{t}^{rs}=Z_{\chi}(
  1+b_{\chi}\mq{rs}+\bar{b}_{\chi}\tr\Mq)
  P_t^{rs}
  \enum
}
便按其夸克内容重正化。

从裸场
\equation{
  A_{\mu}^{rs}(x)=\psibar_r(x)\dirac{\mu}\dirac{5}\psi_s(x),
  \enum
  \nexteq{2.5ex}
  \tilde{A}^{rs}_{\mu}(x)=
  \lambdabar_r(0,x)\dirac{\mu}\dirac{5}\psi_s(x)
  +\psibar_r(x)\dirac{\mu}\dirac{5}\lambda_s(0,x),
  \enum
  \nexteq{2.5ex}
  P^{rs}(x)=\psibar_r(x)\dirac{5}\psi_s(x),
  \enum
  \nexteq{2.5ex}
  \tilde{P}^{rs}(x)=
  \lambdabar_r(0,x)\dirac{5}\psi_s(x)
  +\psibar_r(x)\dirac{5}\lambda_s(0,x),
  \enum
}
出发，改进后的味非单态轴矢流与密度由
\equation{
  \imps{A}{\mu}^{rs}=A^{rs}_{\mu}
  +\ca\ring{\partial}_{\mu}P^{rs}
  +\cta\tilde{A}^{rs}_{\mu},
  \enum
  \nexteq{2.5ex}
  \imp{P}^{rs}=P^{rs}+\ctp\tilde{P}^{rs},
  \enum
}
给出，并且它们按文献~[\cite{SFimp},\cite{ImpNonD}] 重正化：
\equation{
  \rens{A}{\mu}^{rs}=\ZA(1+b_A\mq{rs}+\bar{b}_A\tr\Mq)
  \imps{A}{\mu}^{rs},
  \enum
  \nexteq{2.5ex}
  \ren{P}^{rs}=\ZP(1+b_P\mq{rs}+\bar{b}_P\tr\Mq)
  \imp{P}^{rs}.
  \enum
}
新的改善系数 $\cta,\ctp$
是使涉及非零流时间费米子场的关联函数达到 O($a$) 改进所需要的，
而所有其他系数已经在四维理论中出现
[\cite{SW}--\cite{ImpNonD}]。
在微扰论树图级，
\equation{
  \cfl=\frac{1}{2},
  \qquad
  \cta=\ctp=-\frac{1}{2},
  \qquad
  b_{\chi}=1.
  \enum
}
顺便指出，所有系数 $\bar{b}_X$ 都是 $g_0^4$ 量级
[\cite{ImpNonD}]。


\lussubsec 6.5 场 $\Ptax^{rs}$ 的改善与重正化

在第 4 节导出的手征对称性关系中，
场 $\Ptax^{rs}$ 扮演着突出角色。
这一情形的改善与重正化稍显复杂，
因此与其他场分开讨论。

如 4.2 小节已指出的，
该场只能与至少包含一个拉格朗日乘子场的
局域复合场混合。
电荷共轭、格点对称性与味道对称性随即意味着
该场是乘法可重正化的。
然而有几个场可以在 $a$ 阶与该密度混合，
其中包括两个此前未曾出现的场：
\equation{
  \hat{P}^{rs}(x)=\lambdabar_r(0,x)\dirac{5}\lambda_s(0,x),
  \enum
  \nexteq{2.5ex}
  Q^{rs}(x)=\lambdabar_r(0,x)\dirac{5}\psi_s(x)-
  \psibar_r(x)\dirac{5}\lambda_s(0,x),
  \enum
}
考察表明，改进且重正化后的密度的一种可能选取是
\equation{
  \imp{\Ptax}^{rs}=
  \Ptax^{rs}+\hat{c}_A
  \ring{\partial}_{\mu}
  \tilde{A}_{\mu}^{rs}+\hat{c}_P\hat{P}^{rs},
  \enum
  \nexteq{2.5ex}
  \ren{\Ptax}^{rs}=
  \tilde{Z}_P\bigl\{
  (1+b_{\Ptax}\mq{rs}+\bar{b}_{\Ptax}\tr\Mq)
  \imp{\Ptax}^{rs}
  +\hat{b}_{\Ptax}(\mq{r}-\mq{s})Q^{rs}\bigr\},
  \enum
}
其中照例利用场方程减少了项数。
这些等式右端的场凭借其对称性质或拉格朗日乘子场的内容彼此区分。
它们都是乘法可重正化的，
因此在壳上关联函数中正比于
$\tilde{A}_{\mu}^{rs}$、$\hat{P}^{rs}$、$Q^{rs}$
的抵消项的贡献为 $a$ 阶。

注意到如下关系式即可简化式 (6.25)、(6.26) 中
重正化改进密度的表达式：由于费米子场 Wick 收缩的形式，
作为其直接后果，
\equation{
  \sum_x\langle\Ptax^{rs}(x)P^{sr}_t(y)\rangle
  =\langle S_t^{rr}(y)\rangle+\langle S_t^{ss}(y)\rangle
  +2\cfl\sum_x\langle\hat{P}^{rs}(x)P^{sr}_t(y)\rangle,
  \enum
  \nexteq{2.5ex}
  \sum_x\langle Q^{rs}(x)P^{sr}_t(y)\rangle
  =\langle S_t^{ss}(y)\rangle-\langle S_t^{rr}(y)\rangle,
  \enum
}
对任何 $t>0$ 精确成立。
另一方面，如 4.2 小节所述，
重正化常数可以（而且应当）选得使归一化条件 (4.15)
在连续极限下成立。
由于式 (4.15) 中的关联函数以正比于 $a^2$ 的速率收敛到连续极限，
将其与未重正化的恒等式即式 (6.27)、(6.28) 比较，
便得到
\equation{
  \tilde{Z}_P=1,
  \enum
  \nexteq{2.5ex}
  \hat{c}_P=-2\cfl,
  \quad
  b_{\Ptax}=\bar{b}_{\Ptax}=0,
  \quad
  \hat{b}_{\Ptax}=-\frac{1}{2}b_{\chi}.
  \enum
}
于是重正化密度
\equation{
  \ren{\Ptax}^{rs}=
  \Ptax^{rs}+
  \hat{c}_A
  \ring{\partial}_{\mu}
  \tilde{A}_{\mu}^{rs}-2\cfl\hat{P}^{rs}
  -b_{\chi}\frac{1}{2}(\mq{r}-\mq{s})Q^{rs}
  \enum
}
取相当简单的形式，
其中 $\hat{c}_A$ 是唯一新增的改善系数。


\lussubsec 6.6 格点上的 PCAC 关系

在连续极限下，
只要恰当地选取重正化常数 $Z_A,Z_P$
与重正化夸克质量 $\mR{r}$，
重正化的 PCAC 关系 (4.14) 即成立。
如果再恰当地调节所有改善系数，
则这意味着
\equation{
  \langle\partial\ren{A}^{rs}(x)\rens{P}{t}^{sr}(y)\rangle
  \noenum
  \nexteq{2.0ex}
  \qquad
  =(\mR{r}+\mR{s})\langle\ren{P}^{rs}(x)\rens{P}{t}^{sr}(y)\rangle
  -\langle\ren{\Ptax}^{rs}(x)\rens{P}{t}^{sr}(y)\rangle
  +\rmO(a^2),
  \enum
}
其中依照传统，改进轴矢流的散度取为
\equation{
  \partial\ren{A}^{rs}=\ZA\left(1+b_A\mq{rs}+\bar{b}_A\tr\Mq\right)
  \partial\imp{A}^{rs},
  \enum
  \nexteq{2.5ex}
  \partial\imp{A}^{rs}=
  \ring{\partial}_{\mu}\{A^{rs}_{\mu}+
  \cta\tilde{A}^{rs}_{\mu}\}
  +\ca\drvstar{\mu}\drv{\mu}P^{rs}.
  \enum
}
此处值得再次强调的一个技术细节是：
只要把流时间 $t$ 取为一个以物理单位给出的正值，
PCAC 关系 (6.32) 就对所有 $x$ 成立，包括 $x=y$。


\lussec 7. 关联函数的计算

在数值格点 QCD 中，
涉及非零流时间处的场的关联函数，
可以按照处理普通强子关联函数时通常采取的步骤来计算。
然而流方程必须在某个环节被求解，
这一事实代表了一种需要仔细对待的复杂化。
作为示例，本节针对两个代表性情形把细节具体展开。


\lussubsec 7.1 赝标两点函数

实践中人们关心的是零空间动量下的关联函数
\equation{
  \langle \Pax^{rs}(x)\Pax_t^{sr}(y)\rangle=
  -\sum_{v,w}\langle\tr\{
  [K(t,y;0,v)S(v,x)_{ss}]^{\dagger}K(t,y;0,w)S(w,x)_{rr}\}\rangle
  \enum
}
并利用随机源场 [\cite{MichaelPeisa}]
计算其对一个源点格点集合 $\Gamma$ 的平均，
以减小统计误差。

若取 $x$ 为源点，
则求平均相当于作代换
\equation{
  P^{rs}(x)\to {1\over N_{\Gamma}N_s}
  \sum_{k=1}^{N_s}
  (\psibar_r,\eta_k)(\eta_k,\dirac{5}\psi_s),
  \enum
}
其中 $N_{\Gamma}$ 是 $\Gamma$ 中点的数目，
而
\equation{
  \eta_k(x),\quad k=1,\ldots,N_s,
  \enum
}
是 $\Gamma$ 上随机选取的、均值为零且方差为
\equation{
  \langle\eta_k(v)\eta_l(w)^{\dagger}\rangle=
  \cases{\delta_{kl}\delta_{vw} & if $v,w\in\Gamma$,\cr
  \noalign{\vskip1.0ex}
  0 & otherwise.\cr}
  \enum
}
的复旋量场。
在一个 QCD 模拟中，
计算零空间动量下平均后的关联函数
\equation{
  {1\over N_{\Gamma}}\sum_{x\in\Gamma}\sum_{\vec{y}}
  \langle \Pax^{rs}(x)\Pax_t^{sr}(y)\rangle=
  -{1\over N_{\Gamma}N_s}\sum_{k=1}^{N_s}
  \sum_{\vec{y}}\langle\phi_{k,s}(t,y)^{\dagger}\phi_{k,r}(t,y)\rangle,
  \enum
}
于是相当于对有代表性的规范场系综、
对所有 $k=1,\ldots,N_s$ 与 $h=r,s$，
计算函数
\equation{
  \phi_{k,h}(t,y)=\sum_{w,x}
  K(t,y;0,w)S(w,x)_{hh}\eta_k(x)
  \enum
}
对给定的规范场，
计算 $\phi_{k,h}(t,y)$ 分两步进行：
先通过求解质量为 $m_{0,h}$、源场为 $\eta_k(x)$
的格点狄拉克方程算出 $\phi_{k,h}(0,w)$；
随后必须把算得的场在流时间上从 0 演化到 $t$。
沿流时间增大的方向，
流方程在数值上是稳定的，
用龙格--库塔积分器（附录 D）
便容易以可忽略的积分误差得到解。


\lussubsec 7.2 手征凝聚

对随时间演化的夸克凝聚
\equation{
  \langle S_t^{rr}(x)\rangle=
  -\sum_{v,w}\langle
  \tr\{K(t,x;0,v)\left[S(v,w)_{rr}-\cfl\delta_{vw}\right]
  K(t,x;0,w)^{\dagger}\}\rangle,
  \enum
}
而言，同样希望对其位置 $x$ 求平均。
如上引入随机源场后，即可得到表示式
\equation{
  {1\over N_{\Gamma}}\sum_{x\in\Gamma}
  \langle S_t^{rr}(x)\rangle=
  -{1\over N_{\Gamma}}
  \sum_{v,w}\langle
  \xi_k(t;0,v)^{\dagger}\left[S(v,w)_{rr}-\cfl\delta_{vw}\right]
  \xi_k(t;0,w)\rangle
  \enum
}
其中用到场
\equation{
  \xi_k(t;s,w)=\sum_x K(t;x;s,w)^{\dagger}\eta_k(x)
  \enum
}
在流时间 $s=0$ 处的值。

计算这些场需要求解伴随流方程
\equation{
  (\partial_s+\Delta)\xi_k(t;s,w)=0
  \enum
}
从时刻 $s=t$（此时 $\xi_k$ 与 $\eta_k$ 重合）
积到 $s=0$。
此处不应使用直截了当的龙格--库塔积分，
因为拉普拉斯算子 $\Delta$ 是在由 $s$ 时刻梯度流决定的规范场存在下取的，
那样就必须把规范场沿时间反向演化，
即沿不稳定方向演化。
如附录 E 所解释的，
这一困难可以通过一个避免规范场反向积分的分级方案来克服。


\lussec 8. $\mib\ZA$ 与 ${\mib c}_{\bf fl}$ 的计算策略

只要恰当地选取相关的重正化常数与改善系数，
PCAC 关系 (6.32) 就成立。
如下文所解释的，
反过来，通过要求该关系在运动学参数空间的两个点上精确成立，
可以在分别为 $a^2$ 阶与 $a$ 阶精度内确定常数 $\ZA$
与系数 $\cfl$。

按此思路通过数值模拟计算 $\ZA$ 与 $\cfl$，
假定改善系数 $\csw,\ca$ 已经算出。
空间方向应施加周期性边界条件，
除此之外不再作其他假设。
特别地，无需指明时间方向的边界条件，
因为算得的 $\ZA$ 与 $\cfl$ 的值
只在格距效应的水平上依赖于它们。


\lussubsec 8.1 PCAC 关系的积分形式

不幸的是，当把重正化的 PCAC 关系
用未改进裸场的关联函数展开时，
O($a$) 改进会导致项数膨胀。
不过，将等式对围绕 $y_0$ 的时段 $x_0$
以及空间格点的所有 $\vec{x}$ 求和，
可以实现一点简化。
进入该关系的裸关联函数之一便是
\equation{
  \CPP{rs}(t,d)=\sum_{x_0=y_0-d}^{y_0+d}\sum_{\vec{x}}
  \langle\Pax^{rs}(x)\Pax_t^{sr}(y)\rangle,
  \enum
}
另外还有 $6$ 个这样的关联函数，
其中 $P^{rs}$ 被换成
\equation{
  \tilde{P}^{rs},\hat{P}^{rs},Q^{rs},\ring{\partial}_{\mu}A_{\mu}^{rs},
  \ring{\partial}_{\mu}\tilde{A}_{\mu}^{rs}
  \kern0.5em\hbox{or}\kern0.5em
  \drvstar{\mu}\drv{\mu}P^{rs}.
  \enum
}
对此清单中的最后三个场，
对 $x_0$ 的求和按例如
\equation{
  \CAP{rs}(t,d)=
  \frac{1}{2}\sum_{\vec{x}}\bigl\{
  \langle(\left.A^{rs}_0(x)\right|_{x_0=y_0+d+1}+
          \left.A^{rs}_0(x)\right|_{x_0=y_0+d})
  P^{sr}_t(y)\rangle
  \noenum
  \nexteq{1.0ex}
  {\phantom{\CAP{rs}(t,d)=
  \frac{1}{2}\sum_{\vec{x}}}}\kern-0.65em
  -\langle(\left.A^{rs}_0(x)\right|_{x_0=y_0-d-1}+
           \left.A^{rs}_0(x)\right|_{x_0=y_0-d})
  P^{sr}_t(y)\rangle
  \bigr\},
  \enum
}
的方式化归为求和范围边界处各项之和。

为了记号方便，现在引入缩写
\equation{
  \ZAh{rs}=\ZA\left(1+b_A\mq{rs}+\bar{b}_A\tr\Mq\right),
  \enum
  \nexteq{2.5ex}
  \ZPh{rs}=\ZP\left(1+b_P\mq{rs}+\bar{b}_P\tr\Mq\right).
  \enum
}
通过
\equation{
  m_{rs}={\ZPh{rs}\over\ZAh{rs}}(\mR{r}+\mR{s}),
  \enum
}
定义裸流夸克质量之和后，
积分形式的 PCAC 关系变为
\equation{
  \ZAh{rs}\bigl\{\CAP{rs}+
  \tilde{c}_A\CAtP{rs}+c_A\CdPP{rs}
  -m_{rs}\bigl(\CPP{rs}+\tilde{c}_P\CPtP{rs}\bigr)\bigr\}
  \noenum
  \nexteq{2.0ex}
  \qquad+\CPtP{rs}+\hat{c}_A\CAtP{rs}-2\cfl\CPhP{rs}
  -b_{\chi}\frac{1}{2}(\mq{r}-\mq{s})\CQP{rs}=\rmO(a^2).
  \enum
}
这个等式看起来仍然相当复杂，
但应当记住：
它是若干可计算的关联函数之间的线性关系，
必须对所有流时间 $t>0$、所有时间段 $d\geq0$
以及所有夸克味道 $r\neq s$ 成立。
因此未知系数受到该恒等式的强烈限制。

当 $d$ 大于梯度流的光滑范围 $\sqrt{8t}$ 时，
$\CAtP{rs}$ 按高斯型衰减，
相对式 (8.7) 中其余各项迅速变得可忽略。
而且此时花括号中各项之和以及等式中所有其他项
实际上都与 $d$ 无关。
除非另有说明，
下文只考虑这一 $d$ 范围，
并略去正比于 $\CAtP{rs}$ 的两项。


\lussubsec 8.2 上-下味道道

在上、下夸克质量简并的 QCD 中，
上下道的 PCAC 关系取如下形式：
\equation{
  \ZAh{ud}\bigl\{\CAP{ud}+c_A\CdPP{ud}-m_{ud}\CPP{ud}\bigr\}
  -2\cfl\CPhP{ud}
  =-\bigl(1-\ZAh{ud}\tilde{c}_Pm_{ud}\bigr)\CPtP{ud}
  \enum
}
（为简单起见，从此处起省略 O($a^2$) 误差项）。
裸质量组合 $m_{ud}$ 可以照常由流时间为零处的
PCAC 关系确定，并可在此视为已知。
在两个流时间 $t$ 值处计算式 (8.8) 中的关联函数，
便得到关于比值
\equation{
  {\ZAh{ud}\over1-\ZAh{ud}\tilde{c}_Pm_{ud}}
  \quad\hbox{and}\quad
  {\cfl\over1-\ZAh{ud}\tilde{c}_Pm_{ud}}.
  \enum
}
的两个线性方程。
显然，只要所选的流时间处于标度区
并以物理单位固定，
算得的比值就不依赖于它们，
误差分别不超过 $a^2$ 阶与 $a$ 阶。

当夸克质量接近物理值且格距 $a\leq0.1$ fm 时，
因子 $1-\ZAh{ud}\tilde{c}_Pm_{ud}$
与 1 相差不到百分之一。
该修正项的大小可以通过代入已知的 $m_{ud}$
及 $\tilde{c}_{P}$ 的树图值 (6.22) 来估计。
于是 $\ZAh{ud}$ 与 $\cfl$ 可以在这一小修正因子的精度内算出。
注意：通过改进轴矢流的真空到 $\pi$ 介子矩阵元确定的
裸 $\pi$ 介子衰变常数，
应当用依赖质量的因子 $\ZAh{ud}$ 而非 $\ZA$ 重正化。
不过差别很小，
可以用树图值 $b_A=1$ 作出估计。


\lussubsec 8.3 上-奇异味道道

此道中 PCAC 关系
\equation{
  \ZAh{us}\bigl\{\CAP{us}+c_A\CdPP{us}-m_{us}\CPP{us}\bigr\}
  -2\cfl\CPhP{us}-b_{\chi}\frac{1}{2}(\mq{u}-\mq{s})\CQP{us}
  \noenum
  \nexteq{2.0ex}
  \qquad
  =-\bigl(1-\ZAh{us}\tilde{c}_Pm_{us}\bigr)\CPtP{us}
  \enum
}
出现的附加项正比于质量差 $\mq{u}-\mq{s}$ 的平方，
因此实际中很可能可以忽略。
为了能够确定比值 (8.9)（把 $ud$ 换成 $us$），
除此之外还须在三个流时间 $t$ 值处计算式 (8.10)。

对于 $a\leq0.1$ fm 及物理夸克质量，
修正因子 $1-\ZAh{us}\tilde{c}_Pm_{us}$
与 1 的偏离可达百分之一左右。
与上下道所得结果的比较，
即
\equation{
  {\ZAh{us}(1-\ZAh{ud}\tilde{c}_Pm_{ud})
  \over\ZAh{ud}(1-\ZAh{us}\tilde{c}_Pm_{us})}
  =1+b_A\frac{1}{2}(\mq{s}-\mq{d})+\ZA\tilde{c}_P(m_{us}-m_{ud}),
  \enum
}
不幸只能给出系数 $b_A$ 与 $\tilde{c}_P$
的某种线性组合的估计。


\lussec 9. $2+1$ 味 QCD 中的示例计算

以下几段的目的是演示所提出的
计算重正化常数 $\ZA$、改善系数 $\cfl$
以及低能常数 $\Sigma,F$ 诸策略的可行性。
这一检验只需一次模拟运行即可完成，
并且并不试图真正计算这些低能常数，
因为那将需要把数据外推到无限体积、手征与连续极限。


\lussubsec 9.1 模拟细节

下面报告的计算基于最近一次对 $2+1$ 味夸克 QCD 的模拟
[\cite{TMopenQCD}]。
在该运行（文献~[\cite{TMopenQCD}] 中标记为 $I_1$）中，
模拟了一个在时间方向具有开边界条件 [\cite{openQCD}]
的 $64\times32^3$ 格点，
其参数空间中的位置对应于
估计为 $0.090$ fm 的格距
[\cite{AokiEtAlI},\cite{AokiEtAlII}]，
$\pi$ 介子与 K 介子质量分别约为 203 与 520 MeV
[\cite{TMopenQCD}]。
因此格点的物理尺寸使得
所计算的介子质量与矩阵元的有限体积效应预期很小。

对于该模拟使用的 Iwasaki 规范作用量 [\cite{Iwasaki}]
及标准 O($a$) 改进 Wilson--Dirac 算符
[\cite{SW},\cite{SFimp}]，
系数 $\csw,\ca$ 已非微扰地已知
[\cite{AokiEtAlIII},\cite{KanekoEtAl}]。
该运行产生了由 150 个规范场组态构成的代表性系综，
相邻组态相隔 $9.6$ 个分子动力学时间单位。
为了压制残余的自关联效应，
本节所引的统计误差通过如下方式估计：
把基本观测量系列的测量数据分成每 5 个连续测量为一组的 bin，
再对分箱数据应用 jackknife 方法。

% Fig. 1 (converted from epsf/figure1; graphic from plots/wi.eps -> images/wi.pdf)
\topinsert
\includegraphics[width=10.5cm]{images/wi}
\vspace{0.3cm}
\figurecaption{%
PCAC 关系 (8.8) 中出现的关联函数组合，
绘成求和范围 $d$ 的函数。
所示关联函数的线性组合分别是乘以重正化常数 $\ZAh{ud}$ 的（三角）、
乘以改善系数 $\cfl$ 的（菱形）
以及乘以因子 $1-\ZAh{ud}\ctp m_{ud}$ 的（方块）那些组合，
流时间为 $t=3.86$。
当 $d\geq9$ 时，函数 $\CAtP{ud}$（圆）比其他函数小许多个数量级，
可以安全地忽略（参见 8.1 小节）。
所有量均以格点单位给出。
}
\endinsert


\lussubsec 9.2 $\ZAh{rs}$ 与 $\cfl$ 的计算

进入 PCAC 关系 (8.7) 的关联函数
在流时间 $t$（以格点单位计）取值
$2.47,\,3.86,\,5.56,\,7.56$ 处被计算，
对应的梯度流光滑范围分别为 $0.4,\,0.5,\,0.6,\,0.7$ fm。
为了避免靠近格点在时间方向的边界，
随时间演化场 $P_t^{rs}(y)$ 被插在格点中部。
随后利用时刻 $y_0$ 处的 12 个随机源场，
对各关联函数作空间平移平均。

PCAC 关系中的全部关联函数
都能在整个求和窗口 $d$ 的范围内以很小的统计误差得到。
正如所预期的，
当 $d$ 大于梯度流光滑范围的 2 或 3 倍时，
乘以轴矢流重正化常数的各关联函数之和
与其他关联函数都达到平台
（代表性情形见图~1；
乘以 $\ZAh{ud}$ 的函数在接近格点边界处的轻微下弯
是一种格距效应）。

按照第 8 节概述的策略，
现在通过要求 PCAC 关系 (8.8)
在 $t$ 的两个值 $t_1,t_2$ 处精确成立来确定比值 (8.9)。
只要处于各关联函数的平台范围内，
求和窗口 $d$ 的选取对算得的比值影响很小。
取 $d=18$ 便得到表~1 第二、三列所引的结果，
其中因子 $1-\ZAh{ud}\ctp m_{ud}$
用系数 $\ctp$ 的树图值 (6.22) 作了估计。
不出所料，这一修正非常小（约 $0.2\%$）。

在上-奇异味道道中有两个 O($a$) 质量修正来源。
照上述步骤处理，
来自因子 $1-\ZAh{us}\ctp m_{us}$ 的修正估计约为 $1.3\%$，
而来自式 (8.10) 中正比于 $b_{\chi}$ 项的修正
使结果改变约 $0.2\%$。
如表~1 所示，
两个味道道所得 $\cfl$ 的值之间有显著的一致性。
对流时间的依赖也很小，
而重正化常数 $\ZAh{ud}$ 与 $\ZAh{us}$ 的差别
很可能完全可以由式 (8.4) 中的依赖质量的因子解释。
O($a$) 改进于是进行得非常好，
就重正化常数而言残余格距效应仅为百分之一的几分之一水平。

% Table 1 (re-typeset as LaTeX tabular; data verbatim from source table1)
\topinsert
	ablecaption{$\ZAh{rs}$ 与 $\cfl$ 的模拟结果}
\begin{center}
\renewcommand{\arraystretch}{1.25}
\begin{tabular}{ccccc}
\hline\hline
$t_1,t_2$ & $\ZAh{ud}$ & $\cfl$ & $\ZAh{us}$ & $\cfl$ \\
\hline
$2.47,3.86$ & $0.7929(87)$ & $0.5703(11)$  & $0.8096(41)$ & $0.56867(60)$ \\
$2.47,5.56$ & $0.7904(85)$ & $0.5712(14)$  & $0.8073(40)$ & $0.56990(72)$ \\
$2.47,7.56$ & $0.7895(85)$ & $0.5716(16)$  & $0.8067(39)$ & $0.57025(83)$ \\
$3.86,5.56$ & $0.7885(85)$ & $0.5734(22)$  & $0.8054(40)$ & $0.5726(10)$  \\
$3.86,7.56$ & $0.7883(85)$ & $0.5735(26)$  & $0.8055(39)$ & $0.5724(12)$  \\
$5.56,7.56$ & $0.7882(86)$ & $0.5738(36)$  & $0.8056(38)$ & $0.5723(16)$  \\
\hline\hline
\multicolumn{5}{l}{\footnotesize The values of $\cfl$ in the 3rd and 5th column}\\[-3pt]
\multicolumn{5}{l}{\footnotesize were obtained in the $ud$ and $us$ channel, respectively}\\
\end{tabular}
\end{center}
\endinsert


表~1 所引的重正化常数 $\ZAh{ud}$ 的结果
与文献~[\cite{AokiEtAlIV}] 早先使用基于薛定谔泛函的方法
[\cite{ZASFI},\cite{ZASFII}]
得到的值 $\ZA=0.781(20)$ 一致。
注意改善系数 $\cfl$ 结果上接近其树图值 (6.22)，
尽管在所考虑的规范耦合处并没有很好的理由必然如此。


\lussubsec 9.3 手征凝聚的计算

确定了改善系数 $\cfl$ 之后，
未重正化的随时间演化凝聚 $\Sigma^{uu}_t$
便可按 7.2 小节所述直截了当地计算。
源时刻 $x_0$ 取在格点中部，
以尽量减小激发态对标量密度期望值的贡献。
使用 12 个随机源场并在 $t_1,t_2=3.86,5.56$ 处取 $\cfl$，
计算给出表~2 第二列所引的结果。
在所考虑的流时间范围内，
可见随时间演化的凝聚随 $t$ 增大而减小，粗略如 $1/t$。
统计误差跟随这一演变，
无论如何都小得令人鼓舞。

% Table 2 (re-typeset as LaTeX tabular; data verbatim from source table2)
\topinsert
	ablecaption{格点单位下 $\Sigma^{uu}_t$ 与 $G^{ud}_t$ 的模拟结果}
\begin{center}
\renewcommand{\arraystretch}{1.25}
\begin{tabular}{ccccc}
\hline\hline
$t$ & $\Sigma^{uu}_t$ & $G^{ud}_t$
    & $\Sigma^{uu}_tG^{ud}/G^{ud}_t$ & $\Sigma^{uu}_t/G^{ud}_t$ \\
\hline
$2.47$ & $0.0006277(95)$ & $0.01484(23)$  & $0.003962(61)$ & $0.04230(85)$ \\
$3.86$ & $0.0004251(58)$ & $0.01028(16)$  & $0.003872(55)$ & $0.04134(81)$ \\
$5.56$ & $0.0002911(36)$ & $0.00720(11)$  & $0.003785(51)$ & $0.04040(78)$ \\
$7.56$ & $0.0002017(23)$ & $0.005092(80)$ & $0.003711(48)$ & $0.03962(76)$ \\
\hline\hline
\end{tabular}
\end{center}
\endinsert


为了能把随时间演化的凝聚与手征极限下的
夸克凝聚 $\Sigma$ 联系起来，
还需要计算真空到 $\pi$ 介子矩阵元
\equation{
  \Gpi=\ZP(1+b_P\mq{ud}+\bar{b}_P\tr\Mq)\kern1pt G^{ud},
  \enum
  \nexteq{3.0ex}
  \Gpit=Z_{\chi}(1+b_{\chi}\mq{ud}+\bar{b}_{\chi}\tr\Mq)\kern1pt G^{ud}_t,
  \enum
}
实际上，由于式 (4.24)、(4.25) 中只出现比值
\equation{
 {\rens{\Sigma}{t}\over\Gpit}=
 {\Sigma^{uu}_t\over G^{ud}_t}
 \enum
}
故只需计算 $\Gpi$ 及未重正化的矩阵元 $G^{ud}_t$。
后者可以在大时间间隔 $|x_0-y_0|$ 处
从赝标关联函数 (7.1) 确定，
其方式与通常从普通赝标两点函数提取流时间为零处的
裸矩阵元 $G^{ud}$ 相同。
把 $x_0$ 取在紧邻格点边界的值处，
在所有情形下都存在一个较宽的 $y_0$ 时间范围，
其中单 $\pi$ 介子中间态占主导，
因而可以容易地确定所需的矩阵元。

与随时间演化的凝聚类似，
矩阵元 $G^{ud}_t$ 是 $t$ 的单调递减函数。
然而表~2 最后两列所列的比值随 $t$ 的变化小得多，
这与如下事实一致：
在手征极限下这些比值对流时间的依赖必须消失。


\lussubsec 9.4 若干说明

\noindent{\itshape 换算到物理单位.}
模拟格点的格距 $a=0.08995(40)$ fm 由 PACS-CS 通过计算
$\Omega$ 重子质量确定 [\cite{AokiEtAlII}]。
PACS-CS 还利用步长标度与薛定谔泛函计算了把裸矩阵元 $G^{ud}$
与 $\MSbar$ 方案下 2 GeV 处的 $\Gpi$ 相联系的重正化常数
$\ZP=0.580(21)$ [\cite{AokiEtAlIV}]\kern1pt\sfn{$\dagger$}{%
此处给出的 $\ZP$ 值需要把文献~[\cite{AokiEtAlIV}]
所引的若干表项组合起来并将误差按平方和合成而得到。}。
例如在 $t=3.86$ 处即得结果
\equation{
  {\rens{\Sigma}{t}\Gpi\over\Gpit}=[287(4)\,\MeV]^3,
  \enum
  \nexteq{2.5ex}
  {\rens{\Sigma}{t}\over\Gpit}=91(2)\,\MeV,
  \enum
}
其中忽略了式 (9.1) 中的 O($a$) 质量修正
（反正式 (9.4) 的误差由 $\ZP$ 的误差主导）。
显然，虽然这些比值可能已接近手征极限下的凝聚 $\Sigma$
与衰变常数 $F$，
但要在可控系统误差下确定 $\Sigma$ 与 $F$，
还需要对一系列格距进行模拟。

\vskip0.5ex
\noindent{\itshape 赝标衰变常数.}
由于轴矢流重正化常数的计算不需要单独的模拟，
重正化的赝标介子衰变常数在每个被模拟的格点上都是直接可得的。
不过根据味道道以及所要达到的精度，
用于确定重正化常数 $\ZAh{rs}$ 的 PCAC 关系中的
依赖质量 O($a$) 修正可能需要非微扰地估计。

\vskip0.5ex
\noindent{\itshape 向连续极限的标度.}
本节报告的计算结果依赖于流时间的选择。
当模拟不同格距的格点时，
应当把这些参数以某个合适的参考标度为单位固定，
例如参考流时间 $t_0$ [\cite{WilsonFlow}]。
只有这样，才能期望所计算的重正化量
以正比于 $a^2$ 的速率收敛到连续极限。


\lussec 10. 结束语

通过向夸克场的推广，
梯度流成为研究 QCD 中手征对称性自发破缺
以及轻赝标介子物理的有力工具。
本文的重点是理论框架及其在广泛使用的格点 QCD 表述之一中的实现。
鉴于流的重正化性质，
可以预期：只要格点理论的局域性与规范不变性得到保证，
格点正规化的选择在连续极限下无关紧要。
显然，这一理论分析同样适用于
具有其他规范群与费米子多重态的场论。

在迄今具体完成的例子中，
梯度流在数值格点 QCD 中的应用
相对于既有的计算策略展现出重要的技术优势。
除了第 1 节已经提到的那些之外，
还存在许多潜在的进一步用途。
例如味单态道尚未被考虑，
而随时间演化的凝聚很可能是非零温度 QCD 中
一个易于获取的手征序参量。
此外，把流应用于共形窗附近的类 QCD 理论，
或许可以带来有趣的定性洞见。

本节第 9 节报告的全部计算是在 CERN 的一台专用 PC 机群上完成的。
作者感谢 CERN 管理层提供所需经费，
并感谢 CERN IT 部门的技术支持。



\lusapp A. 记号约定

\vskip-2.5ex

\lussubsec A.1 规范群

$\SUn$ 的李代数可以等同于全部复的、反厄米、无迹 $N\times N$ 矩阵构成的线性空间。
取这类矩阵的一组基 $T^a$，$a=1,\ldots,N^2-1$，
李代数的一般元素 $X$ 即 $X=X^aT^a$，分量 $X^a$ 为实数
（重复的群指标自动求和）。
假定生成元满足
\equation{
  \tr\{T^aT^b\}=-\frac{1}{2}\delta^{ab},
  \enum
}
除此之外不作进一步规定。

规范场的归一化使得规范变换与协变导数
不显含规范耦合。
$D$ 维洛伦兹指标 $\mu,\nu,\ldots$ 取值 $0$ 到 $D-1$，
成对出现时求和。
时空度规取欧几里得度规。特别地，
对任意动量 $p$ 有 $p^2=p_{\mu}p_{\mu}$ 且 $\delta_{\mu\mu}=D$。


\lussubsec A.2 狄拉克矩阵

狄拉克矩阵 $\dirac{\mu}$ 满足
\equation{
  \diracstar{\mu}{\dagger}=\dirac{\mu},\qquad
  \{\dirac{\mu},\dirac{\nu}\}=2\delta_{\mu\nu}.
  \enum
}
标量积 $\dirac{\mu}p_{\mu}$ 照 usual 缩写为
$\slash{p}$，且
$\sigma_{\mu\nu}={i\over2}\left[\dirac{\mu},\dirac{\nu}\right]$。
在任何接近 4 的维数 $D$ 下，
\equation{
  \dirac{5}=\dirac{0}\dirac{1}\dirac{2}\dirac{3},
  \qquad
  \tr\{1\}=4.
  \enum
}
注意：若按这些约定使用维数正规化，
则轴矢流需要重正化。


\lussubsec A.3 格点理论

本文考虑的格点理论照 usual 建立在间距 $a=1$ 的四维超立方格点上。
夸克场采用 Wilson 表述：
夸克场位于格点座上，
带有与连续理论中夸克场相同的指标。
在规范场 $U(x,\mu)$ 存在下，
作用于夸克场 $\psi(x)$ 的规范协变前向与后向差分算符定义为
\equation{
  \nab{\mu}\psi(x)=U(x,\mu)\psi(x+\hat{\mu})-\psi(x),
  \enum
  \nexteq{2.5ex}
  \nabstar{\mu}\psi(x)=\psi(x)-U(x-\hat{\mu},\mu)^{-1}\psi(x-\hat{\mu}),
  \enum
}
其中 $\hat{\mu}$ 表示 $\mu$ 方向的单位矢量。
在格点上，$\drv{\mu}$ 与 $\drvstar{\mu}$ 表示普通的前向与后向差分算符，
而
\equation{
  \ring{\partial}_{\mu}=\frac{1}{2}(\drv{\mu}+\drvstar{\mu})
  \enum
}
表示它们的对称组合。

微分算符 $\partial^a_{x,\mu}$ 按下式作用于规范场 $U$ 的可微函数 $f(U)$：
\equation{
  \partial^a_{x,\mu}f(U)=
  {\rmd\over\rmd s}f(\rme^{sX}U)\!\left.{{\vphantom{r\over s}}}\right|_{s=0},
  \hskip1.7em
  X(y,\nu)=\cases{T^a & if $(y,\nu)=(x,\mu)$,\cr
                          \noalign{\vskip0.8ex}
                        0  & otherwise.\cr}
  \enum
}
虽然这些算符依赖于生成元 $T^a$ 的选取，
可以证明梯度场
\equation{
  (\partial f)(U,x,\mu)=T^a\partial^a_{x,\mu}f(U)
  \enum
}
与基的选取无关。



\lusapp B. 映射 (5.9) 的性质

映射 (5.9) 可微且满足
\equation{
  \Psu(\rme^{sX})\mathrel{\mathop=_{s\to0}}
  sX+\rmO(s^2)
  \enum
}
对所有属于 $\SUn$ 李代数的矩阵 $X$ 成立。
这些性质意味着存在 $\SUn$ 中单位元的一个开邻域，
$\Psu$ 在其中是一一对应的，
且雅可比矩阵
\equation{
  J(U)^{ab}=-2\kern1pt\tr\kern-1pt\left\{T^a\Psu(T^bU)\right\}
  \enum
}
处处非奇异。

利用简单的范数估计可以证明：
区域
\equation{
  \Dsu=\left\{U\in\SUn\bigm|
  \Re\tr\{1-U\}\leq(16N)^{-1}\right\}
  \enum
}
的内部就是这样一个单位元的邻域。
此外，雅可比行列式 $\det J(U)$ 在 $\Dsu$ 上严格为正。
于是，对该 $\SUn$ 不变测度（在李代数上）
相应的狄拉克 $\delta$ 函数而言，
恒等式
\equation{
  \int_{\Dsu} \rmd U\,\delta\left(\Psu(U)-\Psu(V)\right)f(U)=
  k\det J(V)^{-1}f(V)
  \enum
}
对所有可积函数 $f(U)$ 及所有 $V\in\Dsu$ 成立。
此式中的比例常数 $k$ 为正且可以显式求出，
但本文不需要它的数值。



\lusapp C. Wick 收缩

O($a$) 改进格点理论中基本费米子场的 Wick 收缩，
先就时间步长为 $\eps$ 的离散化流方程情形算出。
这样可以避免时序上的歧义，
而在计算的最后容易取极限 $\eps\to0$。


\lussubsec C.1 费米子作用量与泛函积分

作为开始，回顾一下是有帮助的：
$4+1$ 维 O($a$) 改进格点理论的 total 作用量的费米子部分是
一个边界项
\equation{
  \SF=\sum_x\bigl\{\psibar(x)\Dm\psi(x)
  +\cfl\lambdabar(0,x)\lambda(0,x)\bigr\},
  \enum
}
与体作用量 (5.8) 之和。
式 (C.1) 中的 Wilson--Dirac 算符 $\Dm$
包含 Sheikholeslami--Wohlert 项 [\cite{SW},\cite{SFimp}]
以及带裸质量矩阵 $M_0$ 的质量项。

在体作用量 (5.8) 中，
随时间演化场理解为满足边界条件 (2.6)。
于是泛函积分中被积分的、不受约束的费米子场
（反费米子场类似）可取为
\equation{
  \psi(x),\quad\chi(\eps,x),\ldots,\chi(T,x),\quad
  \lambda(0,x),\ldots,\lambda(T-\eps,x).
  \enum
}
由于作用量对这些场是二次的，
费米子积分由 Wick 定理给出。


\lussubsec C.2 场变换

在下文中，满足
\equation{
  K_{\eps}(t,x;t,y)=\delta_{xy}
  \enum
}
（在范围 $0\leq t\leq T$ 内的所有 $t$ 处）
以及离散化夸克流方程
\equation{
  \left(\partial_t-\Delta\right)K_{\eps}(t,x;s,y)=0
  \enum
}
（若 $0\leq s\leq t<T$）
的核 $K_{\eps}(t,x;s,y)$ 扮演着重要角色。
显然，由于流方程 (C.4) 可以从时刻 $t=s$
开始一步步求解，
这些方程在所有 $t\geq s$ 处唯一确定该核。

通过一个场变换可以简化费米子传播子的计算：
将 $\chi(t,x)$ 通过 $\psi(x)$ 与另一个定义在
$t=\eps,\ldots,T$ 的场 $\phi(t,x)$ 表示为
\equation{
  \chi(t,x)=\sum_yK_{\eps}(t,x;0,y)\psi(y)
  +\eps\sum_{s=\eps}^t\sum_yK_{\eps}(t,x;s,y)\phi(s,y).
  \enum
}
对任意给定的场 $\psi(x)$，
由场变量 $\chi(\eps,x),\ldots,\chi(T,x)$
到新变量 $\phi(\eps,x),\ldots,\phi(T,x)$ 的映射
由于恒等式
\equation{
  \phi(t+\eps,x)=\cases{
  \left(\partial_t-\Delta\right)\chi(t,x) & if $0<t<T$,\cr
  \noalign{\vskip1.5ex}
  {1\over\eps}(\chi(\eps,x)-\psi(x))-\Delta\psi(x) &
  if $t=0$.\cr
  }
  \enum
}
而是可逆的。
而且相应的雅可比矩阵是下三角的，
其行列式不依赖规范场。

经过该场变换以及反费米子场的相应变换后，
费米子作用量取如下形式：
\equation{
  \SF+\SFfl=\sum_x
  \bigl\{\psibar(x)\Dm\psi(x)+\cfl\lambdabar(0,x)\lambda(0,x)\bigr\}
  \noenum
  \nexteq{0.0ex}
  {\phantom{\SF+\SFfl=}}
  +\eps\sum_{t=0}^{T-\eps}\sum_x\bigl\{
  \lambdabar(t,x)\phi(t+\eps,x)+\phibar(t+\eps,x)\lambda(t,x)\bigr\}.
  \enum
}
这一变换从而实现了场变量之间近乎完全的退耦。


\lussubsec C.3 收缩

由于基本夸克场 $\psi,\psibar$ 与其他场完全退耦，
这些场的收缩
\equation{
  \wick{\psi(x)\psibar(y)}=S(x,y),
  \enum
  \nexteq{2.5ex}
  \Dm S(x,y)=\delta_{xy},
  \enum
}
如预期那样与四维理论中的夸克传播子一致。
由作用量 (C.7) 导出的其余非零收缩为
\equation{
  \verylongwick{\phi(t+\eps,x)\lambdabar(t,y)}=
  \longwick{\lambda(t,x)\phibar(t+\eps,y)}={1\over\eps}\delta_{xy},
  \quad t=0,\eps,\ldots T-\eps,
  \enum
  \nexteq{2.5ex}
  \longwick{\phi(\eps,x)\phibar(\eps,y)}=-{\cfl\over\eps^2}\delta_{xy}.
  \enum
}
对于涉及场 $\chi$ 与 $\chibar$ 的收缩则得到
\equation{
  \longwick{\chi(t,x)\lambdabar(s,y)}=
  \theta(t>s)K_{\eps}(t,x;s+\eps,y),
  \enum
  \nexteq{3.0ex}
  \longwick{\lambda(t,x)\chibar(s,y)}=
  \theta(t<s)K_{\eps}(s,y;t+\eps,x)^{\dagger},
  \enum
  \nexteq{3.0ex}
  \longwick{\chi(t,x)\psibar(y)}=\sum_{v}
  K_{\eps}(t,x;0,v)S(v,y),
  \enum
  \nexteq{2.4ex}
  \wick{\psi(x)\chibar(s,y)}=\sum_{w}
  S(x,w)K_{\eps}(s,y;0,w)^{\dagger},
  \enum
  \nexteq{2.4ex}
  \longwick{\chi(t,x)\chibar(s,y)}=\sum_{v,w}
  K_{\eps}(t,x;0,v)S(v,w)
  K_{\eps}(s,y;0,w)^{\dagger}
  \noenum
  \nexteq{1.5ex}
  {\phantom{\longwick{\chi(t,x)\chibar(s,y)}=}}
  -\cfl\sum_vK_{\eps}(t,x;\eps,v)K_{\eps}(s,y;\eps,v)^{\dagger}.
  \enum  
}
在所有这些等式中，
$\chi$ 与 $\chibar$ 的流时间变量取值于 $(0,T]$，
而 $\lambda$ 与 $\lambdabar$ 的流时间变量取值于 $[0,T)$。


\lussubsec C.4 连续时间极限下的收缩

当时间步长 $\eps$ 趋于零时，
核 $K_{\eps}(t,x;s,y)$ 光滑地收敛到格点理论连续时间表述下
夸克流方程的基本解 $K(t,x;s,y)$。
因此，由式 (C.8) 与 (C.12)--(C.16) 给出的
费米子场全部非零收缩的集合具有良定义的极限。



\lusapp D. 格点流方程的积分

纯规范梯度流的数值积分已在文献~[\cite{WilsonFlow}] 的附录 C 中讨论过。
这里把该文献中描述的三阶龙格--库塔积分器推广到夸克流方程。


\lussubsec D.1 梯度流的积分

仿照文献~[\cite{WilsonFlow}]，
把流方程 (5.2) 写成抽象形式
\equation{
   \partial_tV_t=Z(V_t)V_t,
   \enum
}
其中规范场 $V_t$
被视为一个（高维）李群的元素，
而 Wilson 作用量的梯度 $Z(V_t)$ 是相应李代数的元素。
上述龙格--库塔积分器以步长 $\eps$ 推进。
设 $V_t$ 在某流时间 $t$ 处已知，
通过计算场
\equation{
  W_0=V_t,
  \noenum
  \nexteq{2.5ex}
  W_1=\exp\bigl\{\frac{1}{4}Z_0\bigr\}W_0,
  \noenum
  \nexteq{2.5ex}
  W_2=\exp\bigl\{\frac{8}{9}Z_1-\frac{17}{36}Z_0\bigr\}W_1,
  \noenum
  \nexteq{2.5ex}
  W_3=\exp\bigl\{\frac{3}{4}Z_2
                 -\frac{8}{9}Z_1+\frac{17}{36}Z_0\bigr\}W_2,
  \enum
}
其中
\equation{
  Z_i=\eps Z(W_i),\qquad i=0,1,2.
  \enum
}
便得到时刻 $t+\eps$ 处精确解 $V_{t+\eps}$ 的近似。
这些规则使得
\equation{
  W_3=V_{t+\eps}+\rmO(\eps^4)
  \enum
}
因而 $W_3$ 给出所期望的 $V_{t+\eps}$ 的近似。


\lussubsec D.2 夸克流方程的积分

同样有益的是把流方程改写成抽象形式
\equation{
  \partial_t\chi_t=\Delta(V_t)\chi_t,
  \enum
}
其中夸克场 $\chi_t$ 被视为一个复矢量空间的元素。
设在某流时间 $t$ 处给定了 $V_t$ 与 $\chi_t$，
并假定规范场 $W_0,W_1,W_2$ 如上，
则夸克场
\equation{
  \phi_0=\chi_t,
  \noenum
  \nexteq{2.5ex}
  \phi_1=\phi_0+\frac{1}{4}\Delta_0\phi_0,
  \noenum
  \nexteq{2.5ex}
  \phi_2=\phi_0+\frac{8}{9}\Delta_1\phi_1-\frac{2}{9}\Delta_0\phi_0,
  \noenum
  \nexteq{2.5ex}
  \phi_3=\phi_1+\frac{3}{4}\Delta_2\phi_2,
  \enum
}
可以如下计算：
\equation{
  \Delta_i=\eps\Delta(W_i),\qquad i=0,1,2.
  \enum
}
不难证明
\equation{
  \phi_3=\chi_{t+\eps}+\rmO(\eps^4),
  \enum
}
于是 $\phi_3$ 对 $\chi_{t+\eps}$ 的逼近精度
与纯规范流积分器所达到的精度相匹配。



\lusapp E. 伴随流方程 (7.10) 的积分

伴随流方程的数值积分之所以复杂，
是因为规范场流方程的反向积分是指数不稳定的。
本附录先把问题重新表述，
然后讨论其解法。
为简单起见，略去式 (7.10) 中的指标 $k$。


\lussubsec E.1 伴随龙格--库塔积分

对给定步长 $\eps>0$ 及 $t=n\eps$，$n=0,1,2,\ldots$，
令
\equation{
  V_t^{\eps}(x,\mu)\quad\hbox{and}\quad
  \chi_t^{\eps}(x)
  \enum
}
为由附录 D 定义的龙格--库塔积分器生成的规范场与夸克场。
存在另一列夸克场
\equation{
  \xi_s^{\eps}(x),\quad s=t,t-\eps,t-2\eps,\ldots,0,
  \enum
}
其初值为
\equation{
  \xi_t^{\eps}(x)=\eta(x),
  \enum
}
且恒等式
\equation{
  (\xi_s^{\eps},\chi_s^{\eps})=(\eta,\chi_t^{\eps})
  \enum
}
对所有 $s$ 精确成立。
这一性质保证 $\xi_s^{\eps}$ 是伴随流方程的一个近似解，
精度达 $\eps^3$ 阶，
并具有所需的初值。

从 $\chi_s^{\eps}$ 推进到 $\chi_{s+\eps}^{\eps}$
的龙格--库塔步骤相当于把算符
\equation{
  R_s^{\eps}=1+\frac{1}{4}\Delta_0+\frac{3}{4}\Delta_2
  \bigl\{1-\frac{2}{9}\Delta_0+\frac{8}{9}\Delta_1(1+\frac{1}{4}\Delta_0)
  \bigr\}
  \enum
}
作用于 $\chi_s^{\eps}$，
其中 $\Delta_i$ 由式 (D.7) 给出，
而其中出现的场 $W_i$ 由取 $W_0=V_s^{\eps}$ 的式 (D.2) 给出。
由式 (E.4) 即可推知
\equation{
  \xi_s^{\eps}(x)=({R_s^{\eps}}^{\dagger}\xi_{s+\eps}^{\eps})(x)
  \enum
  \nexteq{2.5ex}
  {R_s^{\eps}}^{\dagger}=1+\frac{1}{4}\Delta_0+
  \bigl\{1-\frac{2}{9}\Delta_0+(1+\frac{1}{4}\Delta_0)\frac{8}{9}\Delta_1
  \bigr\}\frac{3}{4}\Delta_2,
  \enum
}
即 $\xi_s^{\eps}$ 通过计算
\equation{
  \lambda_3=\xi_{s+\eps}^{\eps},
  \noenum
  \nexteq{2.5ex}
  \lambda_2=\frac{3}{4}\Delta_2\lambda_3,
  \noenum
  \nexteq{2.5ex}
  \lambda_1=\lambda_3+\frac{8}{9}\Delta_1\lambda_2,
  \noenum
  \nexteq{2.5ex}
  \lambda_0=\lambda_1+\lambda_2+
  \frac{1}{4}\Delta_0(\lambda_1-\frac{8}{9}\lambda_2),
  \enum
}
并由 $\xi_s^{\eps}=\lambda_0$ 得到。


\lussubsec E.2 稳定性问题

伴随递推 (E.8) 是一个光滑化操作，因而在数值上是稳定的。
然而为了能在时刻 $s+\eps$ 应用该递推，
必须已知流时间 $s$ 处的规范场 $V_s^{\eps}$。
原则上，从 $s=t$ 到 $s=0$ 的整列规范场
可以通过梯度流的反向积分重构，
但这样的策略并不实用，
因为沿此方向流是指数不稳定的。

另一种办法是从基本场 $V_0^{\eps}=U$ 出发
通过流的正向积分计算 $V_s^{\eps}$。
若对每个 $s$ 都重新计算一遍，
整个积分所需的规范场更新步数
将按时间步数 $m=t/\eps$ 的平方增长。
该方法在数值上是安全的，
算得的夸克场实际上以机器精度满足恒等式 (E.4)。
但其计算代价可能大得令人却步。


\lussubsec E.3 改进方案

把流时间序列 $0,\eps,\ldots,t$ 分成 $n_b$ 个大小相近的连续区块，
可以使伴随流方程的积分更有效率。
做法是：先用正向积分从时刻 0 算出最后一个区块起点时刻 $r$ 处的场
$V_r^{\eps}$ 并将其保存于计算机内存；
随后从时刻 $t$ 到 $r$ 积分伴随流方程，
其间应用更新步骤 (E.8) 并使用上述安全的规范场重构方法，
只是现在重构从存储的场而非时刻 0 的场开始。

下一步，计算倒数第二个区块中第一个时刻 $r$ 处的规范场并存入内存，
然后伴随流方程的积分照前述方式继续推进到时刻 $r$。
显然各区块可以如此依次处理，
直至积分到达时刻 $0$。

当 $m=n_b^2$ 时该方案的规范场更新总步数最少。
虽然 $m$ 很少恰为整数的平方，
但总可以取 $n_b=\lceil m^{1/2}\rceil$
并把 $m$ 分成大小为 $m_b=\lfloor m/n_b\rfloor$ 与 $m_b+1$ 的区块。
此时计算量近似按 $m^{3/2}$ 标度，
而且在很小的 $m$ 值处就已远小于上一小节安全方案所需的计算量。


\lussubsec E.4 分级方案

进一步加速的办法是把每个区块再分成 $n_b$ 个更小的区块，
后者再分成 $n_b$ 个更小的区块，依此类推，
并在每个区块层级保存一个规范场。
给定 $m$ 与已保存规范场的数目 $n_s$，
接近最优的区块数选取为
\equation{
  n_b=\lceil m^{1/(n_s+1)}\rceil.
  \enum
}
例如对 $m\leq10^3$、$n_s=4$，
需要执行的规范场更新总数不超过约 $5m$。
此外，当多个夸克场同时演化时，
中间规范场只需计算一次。

\begin{thebibliography}{99}



% Wilson flow & emergence of the topological sectors

\bibitem{WilsonFlow}
M. L\"uscher,
{\it Properties and uses of the Wilson flow in lattice QCD},
JHEP 1008 (2010) 071


% Renormalization of the gradient flow

\bibitem{RenFlow}
M. L\"uscher, P. Weisz,
{\it Perturbative analysis of the gradient flow in non-Abelian gauge
theories},
JHEP 1102 (2011) 051


% Step scaling

\bibitem{StepScaling}
M. L\"uscher, P. Weisz, U. Wolff,
{\it A numerical method to compute the running coupling in asymptotically 
free theories},
Nucl. Phys. B359 (1991) 221

\bibitem{NogradiEtAl}
Z. Fodor, K. Holland, J. Kuti, D. Nogradi, C. H. Wong,
{\it The Yang-Mills gradient flow in finite volume},
JHEP 1211 (2012) 007; {\it The gradient flow running coupling scheme},
PoS (Lattice 2012) 050

\bibitem{FritzschRamos}
P. Fritzsch, A. Ramos,
{\it The gradient flow coupling in the Schr\"odinger functional},
arXiv:1301.4388


% Quark source smearing

\bibitem{Guesken}
S. G\"usken,
{\it A study of smearing techniques for hadron correlation functions},
Nucl. Phys. B (Proc. Suppl.) 17 (1990) 361

\bibitem{AlexandrouEtAl}
C. Alexandrou, F. Jegerlehner, S. G\"usken, K. Schilling, R. Sommer,
{\it B meson properties from lattice QCD},
Phys. Lett. B256 (1991) 60


% Renormalization of the Langevin equation

\bibitem{Zinn}
J. Zinn--Justin,
{\it Renormalization and stochastic quantization},
Nucl. Phys. B275 [FS17] (1986) 135

\bibitem{ZinnZwanziger}
J. Zinn--Justin, D. Zwanziger, 
{\it Ward identities for the stochastic quantization of gauge fields},
Nucl. Phys. B295 [FS21] (1988) 297


% Renormalization of quantum field theories on a half space

\bibitem{Symanzik}
K. Symanzik,
{\it Schr\"odinger representation and Casimir effect in
renormalizable quantum field theory},
Nucl. Phys. B190 [FS3] (1981) 1


% Wilson formulation of lattice QCD

\bibitem{Wilson}
K. G. Wilson, {\it Confinement of quarks}, Phys. Rev. D10 (1974) 2445


% Exact chiral symmetry on the lattice

% Ginsparg-Wilson relation

\bibitem{GinspargWilson}
P. H. Ginsparg, K. G. Wilson,
{\it A remnant of chiral symmetry on the lattice},
Phys. Rev. D25 (1982) 2649


% Domain wall fermions

\bibitem{Kaplan}
D. B. Kaplan,
{\it A method for simulating chiral fermions on the lattice},
Phys. Lett. B288 (1992) 342;
{\it Chiral fermions on the lattice},
Nucl. Phys. B (Proc. Suppl.) 30 (1993) 597

\bibitem{Shamir}
Y. Shamir,
{\it Chiral fermions from lattice boundaries},
Nucl. Phys. B406 (1993) 90

\bibitem{FurmanShamir}
V. Furman, Y. Shamir,
{\it Axial symmetries in lattice QCD with Kaplan fermions},
Nucl. Phys. B439 (1995) 54


% Classically perfect Dirac operator

\bibitem{Hasenfratz}
P. Hasenfratz, {\it Prospects for perfect actions},
Nucl. Phys. B (Proc. Suppl.) 63 (1998) 53;
{\it Lattice QCD without tuning, mixing and current renormalization},
Nucl. Phys. B525 (1998) 401

\bibitem{HLN}
P. Hasenfratz, V. Laliena, F. Niedermayer,
{\it The index theorem in QCD with a finite cutoff},
Phys. Lett. B427 (1998) 125


% Neuberger-Dirac operator

\bibitem{NeubergerDirac}
H. Neuberger,
{\it Exactly massless quarks on the lattice},
Phys. Lett. B417 (1998) 141;
{\it More about exactly massless quarks on the lattice},
{\it ibid.}\/ B427 (1998) 353


% Lattice chiral symmetry

\bibitem{ExactChSy}
M. L\"uscher,
{\it Exact chiral symmetry on the lattice and the Ginsparg--Wilson relation},
Phys. Lett. B428 (1998) 342

\bibitem{BoulderReview}
F. Niedermayer,
{\it Exact chiral symmetry, topological charge and related topics},
Nucl. Phys. B (Proc. Suppl.) 73 (1999) 105


% SU(2) chiral perturbation theory

\bibitem{ChPT}
J. Gasser, H. Leutwyler,
{\it Chiral perturbation theory to one loop},
Ann. Phys. 158 (1984) 142


% Global existence of the smoothing flow

\bibitem{ArnoldI}
V. I. Arnold,
{\it Ordinary Differential Equations}, 3rd ed. (Springer-Verlag, Berlin, 2008)


% Stout smearing

\bibitem{Stout}
C. Morningstar, M. Peardon, 
{\it Analytic smearing of SU(3) link variables in lattice QCD},
Phys. Rev. D69 (2004) 054501


% Symanzik effective action

\bibitem{SymanzikSeffI}
K. Symanzik, 
{\it Some topics in quantum field theory},
in: Mathematical problems in theoretical physics, eds. R. Schrader
et al., Lecture Notes in Physics, Vol. 153 (Springer, New York, 1982)

\bibitem{SymanzikSeffII}
K. Symanzik,
{\it Continuum limit and improved action in lattice theories:
(I). Principles and $\varphi^4$ theory},
Nucl. Phys. B226 (1983) 187


% On-shell improvement

\bibitem{OnShell}
M. L\"uscher, P. Weisz,
{\it On-shell improved lattice gauge theories},
Commun. Math. Phys. 97 (1985) 59 [E: {\it ibid.} 98 (1985) 433]


% O(a) improvement

\bibitem{SW}
B. Sheikholeslami, R. Wohlert, 
{\it Improved continuum limit lattice action for QCD with Wilson fermions},
Nucl. Phys. B259 (1985) 572

\bibitem{SFimp}
M. L\"uscher, S. Sint, R. Sommer, P. Weisz,
{\it Chiral symmetry and O(a) improvement in lattice QCD},
Nucl. Phys. B478 (1996) 365


% O(a) improvement with non-degenerate quarks

\bibitem{ImpNonD}
T. Bhattacharya, R. Gupta, W. Lee, S. R. Sharpe, J. M. S. Wu,
{\it Improved bilinears in lattice QCD with nondegenerate quarks},
Phys. Rev. D73 (2006) 034504


% Random sources

\bibitem{MichaelPeisa}
C. Michael, J. Peisa,
{\it Maximal variance reduction for stochastic propagators with 
applications to the static quark spectrum},
Phys. Rev. D58 (1998) 034506


% QCD with open boundary conditions and tm reweighting

\bibitem{TMopenQCD}
M. L\"uscher, S. Schaefer,
{\it  	
Lattice QCD with open boundary conditions and twisted-mass reweighting},
Comput. Phys. Commun. 184 (2013) 519


% Open boundary conditions

\bibitem{openQCD}
M. L\"uscher, S. Schaefer,
{\it Lattice QCD without topology barriers},
JHEP 1107 (2011) 036


% PACS-CS physical point simulations

\bibitem{AokiEtAlI}
S. Aoki et al.~(PACS-CS collab.), 
{\it 2+1 flavor lattice QCD toward the physical point},
Phys. Rev. D79 (2009) 034503

\bibitem{AokiEtAlII}
S. Aoki et al.~(PACS-CS collab.),
{\it Physical point simulation in 2+1 flavor lattice QCD},
Phys. Rev. D81 (2010) 074503


% Iwasaki action

\bibitem{Iwasaki}
Y. Iwasaki,
{\it Renormalization group analysis of lattice theories and
improved lattice action. II -- Four-dimensional non-Abelian SU(N) gauge model},
preprint UTHEP-118 (1983) and arXiv:1111.7054


% Non-perturbative O(a) improvement

\bibitem{AokiEtAlIII}
S. Aoki et al. (CP-PACS collab.),
{\it Nonperturbative O(a) improvement of the Wilson quark action with the 
RG-improved gauge action using the Schr\"odinger functional method},
Phys. Rev. D73 (2006) 034501

\bibitem{KanekoEtAl}
T. Kaneko et al. (CP-PACS/JLQCD and ALPHA collab.),
{\it Non-perturbative improvement of the axial current with 
three dynamical flavors and the Iwasaki gauge action},
JHEP 0704 (2007) 092


% Computation of ZA and ZP

\bibitem{AokiEtAlIV}
S. Aoki et al. (PACS-CS collab.),
{\it Non-perturbative renormalization of quark mass 
in $N_f$=2+1 QCD with the Schr\"odinger functional scheme},
JHEP 1008 (2010) 101


% SF method for ZA

\bibitem{ZASFI}
M. L\"uscher, S. Sint, R. Sommer, H. Wittig,
{\it Non-perturbative determination of the axial current 
normalization constant in O(a) improved lattice QCD},
Nucl. Phys. B491 (1997) 344

\bibitem{ZASFII}
M. Della Morte, R. Hoffmann, F. Knechtli, R. Sommer, U. Wolff,
{\it Non-pertur\-ba\-tive renormalization of the axial current 
with dynamical Wilson fermions},
JHEP 0507 (2005) 007





\end{thebibliography}


\end{document}
